\documentclass[11pt]{article}
\usepackage[margin=1in]{geometry}
\usepackage{amsmath,amssymb,bm,mathtools}
\usepackage[round,authoryear]{natbib}
\usepackage{hyperref}
\usepackage{enumitem}
\usepackage{booktabs}
\usepackage{listings}
\usepackage{xcolor}
\usepackage{graphicx}

\hypersetup{colorlinks=true,linkcolor=blue,urlcolor=blue,citecolor=blue}

\lstset{
  basicstyle=\ttfamily\small,
  breaklines=true,
  frame=single,
  columns=fullflexible,
  backgroundcolor=\color{gray!5}
}

\newcommand{\E}{\mathbb E}
\newcommand{\Var}{\operatorname{Var}}
\newcommand{\Cov}{\operatorname{Cov}}
\newcommand{\diag}{\operatorname{diag}}
\newcommand{\tr}{\operatorname{tr}}
\newcommand{\argmin}{\operatorname*{arg\,min}}

\title{The Simultaneous Kalman Smoother:\\
State Space Models Made Easy}
\author{%
Mico Mrkaic\\
International Monetary Fund\\
\texttt{mmrkaic@imf.org}%
\thanks{%
This note and the accompanying \texttt{simksmoother} software package
(Octave/MATLAB and Python) were produced under the author's guidance
and direction, with the assistance of AI (Claude, Anthropic). The
author is responsible for all content, errors, and omissions.}%
}

\date{June 2026}

\begin{document}
\maketitle

\begin{abstract}
The Kalman filter and the Rauch--Tung--Striebel smoother are usually
derived through a forward--backward recursion of conditional means and
covariances.  For economists this derivation is needlessly hard to
follow.  The recursive form has two original motivations, neither of
which now applies to us.  First, it was developed for online vehicle
guidance on 1960s computers with kilobytes of memory, where storing the
joint distribution of an entire latent path was impossible.  Second, it
is in the idiom of control engineering --- state machines, predictions,
online updates --- whereas econometrics is built around joint
distributions, regressions, and batch data.  This note derives the
smoother in the idiom that is native to us.  The key observation is
that stacking the full latent path and all observations produces one
big multivariate normal vector, and the entire Kalman smoother is just
the textbook conditional-Gaussian formula applied to that vector.
Solving the resulting system in precision form is one sparse
block-tridiagonal linear equation $J\widehat X=h$.  The Kalman filter
and the RTS smoother are then exposed as the block-LDL$^\top$ forward
sweep and the selected-inversion backward sweep of the same system ---
implementation details, not theory.  Posterior covariances, marginal
likelihood, missing data, diffuse priors, EM, the Hodrick--Prescott
filter, and exogenous inputs all become small edits to $J$ and $h$.
Heavy-tailed disturbances (Student-$t$, Laplace, Huber) become
iteratively reweighted versions of the same sparse solve, and the
framework extends cleanly to regime-switching dynamics (variational
EM on the same $J$) and to inequality-constrained smoothing
(interior-point QP on the same $J$).  The precision-based view is not
new; it appears in spatial statistics and Bayesian econometrics as
the ``Gaussian Markov random field'' view of state space.  It is,
however, considerably easier to teach.
\end{abstract}

\section{An Example to Motivate the Discussion}
\label{sec:motivation}

Most economists meet the Hodrick--Prescott filter as a recipe: to
extract a trend $\tau_t$ from a series $y_t$, solve
\begin{equation}
\min_{\tau_1,\dots,\tau_T}\ \sum_{t=1}^{T}(y_t-\tau_t)^2
   + \lambda \sum_{t=2}^{T-1}\bigl[(\tau_{t+1}-\tau_t)-(\tau_t-\tau_{t-1})\bigr]^2,
\label{eq:hp-recipe}
\end{equation}
with $\lambda=1600$ for quarterly data, and read off $\widehat\tau$.
The first sum keeps the trend close to the data; the second penalizes
changes in the slope of the trend, with $\lambda$ governing how hard.
Most users stop there.  But where does this objective come from, what
model does it correspond to, and what is $\lambda$, really?  Answering
those questions is the cleanest possible introduction to everything in
this note, because the HP filter turns out to be a one-line state-space
model, and \eqref{eq:hp-recipe} turns out to be a conditional
expectation in a multivariate normal.  Once that is clear, a whole
catalogue of extensions opens up almost for free.

\subsection{The HP filter is a state-space model}

Write potential output $\tau_t$ as a latent process whose
\emph{second difference} is white noise, and observed output $y_t$ as
potential plus a transitory gap:
\begin{align}
(\tau_{t+1}-\tau_t)-(\tau_t-\tau_{t-1}) &= e_t, & e_t &\sim N(0,\sigma_e^2),
\label{eq:hp-state}\\
y_t &= \tau_t + \varepsilon_t, & \varepsilon_t &\sim N(0,\sigma_\varepsilon^2).
\label{eq:hp-meas}
\end{align}
Equation~\eqref{eq:hp-state} is the economic content of the HP filter,
usually left implicit: \emph{potential output is an $I(2)$ process}.
Its level is non-stationary, its growth rate $g_t\equiv\tau_t-\tau_{t-1}$
is non-stationary, but the \emph{change} in its growth rate is
unforecastable noise.  Equation~\eqref{eq:hp-meas} is the trend--cycle
decomposition: the part of output we observe is potential plus a
transitory deviation $\varepsilon_t$, the ``output gap'' up to sign.

This is already a useful step, independent of anything that follows.
The recipe \eqref{eq:hp-recipe} hides a strong, testable economic
assumption --- that potential output is $I(2)$ --- behind a smoothing
objective.  Writing the model in the form \eqref{eq:hp-state}--\eqref{eq:hp-meas}
makes the assumption explicit, and immediately invites the question of
whether it is the right one for a given series.  If potential growth
is better described as \emph{stationary} ($I(1)$ in the level) rather
than as a driftless random walk ($I(2)$ in the level), the HP filter is
misspecified, and one would want a different state equation.  One cannot
even ask that question from the recipe alone.

\subsection{Where \texorpdfstring{$\lambda$}{lambda} comes from}

The two innovation variances are not separately identified by the
\emph{shape} of the extracted trend --- only their ratio is.  To see
this, write the joint Gaussian log-density of the latent path and the
data implied by \eqref{eq:hp-state}--\eqref{eq:hp-meas}.  Dropping
constants, minimizing $-2\log p$ over $\tau$ gives
\begin{equation}
\min_{\tau}\ \frac{1}{\sigma_\varepsilon^2}\sum_{t=1}^{T}(y_t-\tau_t)^2
   + \frac{1}{\sigma_e^2}\sum_{t=2}^{T-1}\bigl[(\tau_{t+1}-\tau_t)-(\tau_t-\tau_{t-1})\bigr]^2.
\label{eq:hp-mle}
\end{equation}
Multiplying through by $\sigma_\varepsilon^2$ leaves the minimizer
unchanged and produces \eqref{eq:hp-recipe} exactly, with
\begin{equation}
\boxed{\ \lambda=\frac{\sigma_\varepsilon^2}{\sigma_e^2}.\ }
\label{eq:lambda}
\end{equation}
So $\lambda$ is not a free smoothing knob: it is the ratio of the
measurement-noise variance to the trend-innovation variance.  The
famous value $\lambda=1600$ for quarterly data is an assertion that the
transitory component of output is $1600$ times as variable as the
shocks to trend growth.  Stated that way, it is a falsifiable claim
about the data, not a tuning choice --- and indeed $\sigma_\varepsilon^2$
and $\sigma_e^2$ can be estimated by maximum likelihood rather than
assumed, which is one of the extensions below.

\subsection{The same problem in simultaneous form}

Now stack.  Collect the trend into $\bm\tau=(\tau_1,\dots,\tau_T)'$ and
the data into $\bm y=(y_1,\dots,y_T)'$.  The measurement
equation~\eqref{eq:hp-meas} stacks to
\begin{equation}
\bm y = \bm\tau + \bm\varepsilon, \qquad \bm\varepsilon\sim N(0,\sigma_\varepsilon^2 I_T).
\label{eq:hp-meas-stacked}
\end{equation}
The state equation~\eqref{eq:hp-state} is a statement about the second
difference of $\bm\tau$.  Let $D_2$ be the $(T-2)\times T$
second-difference operator, the matrix whose $t$-th row places
$(1,-2,1)$ in columns $t,t+1,t+2$:
\begin{equation}
D_2=\begin{pmatrix}
1 & -2 & 1 & & \\
  & 1 & -2 & 1 & \\
  &   & \ddots & \ddots & \ddots \\
\end{pmatrix},
\qquad
D_2\bm\tau = \bm e,\quad \bm e\sim N(0,\sigma_e^2 I_{T-2}).
\label{eq:hp-state-stacked}
\end{equation}
Both \eqref{eq:hp-meas-stacked} and \eqref{eq:hp-state-stacked} are
linear in $\bm\tau$ with Gaussian right-hand sides.  Reading off the
pieces, $\bm\tau$ has prior precision $\sigma_e^{-2}D_2'D_2$ (this is
what the $I(2)$ assumption contributes), and the data contributes
measurement precision $\sigma_\varepsilon^{-2}I_T$.  One qualification
is needed: $D_2'D_2$ has rank $T-2$, so the prior on $\bm\tau$ is an
\emph{intrinsic} (improper) Gaussian --- it is flat in the two
directions spanned by the initial level and slope of the trend, the
diffuse-limit prior on those components.  Strictly, then,
$(\bm\tau,\bm y)$ is jointly Gaussian in the intrinsic/improper sense
rather than a proper multivariate normal.  None of what follows is
affected: the \emph{posterior} of $\bm\tau$ given $\bm y$ is a proper
Gaussian, because the measurement term adds the full-rank precision
$\sigma_\varepsilon^{-2}I_T$.

The trend we want is the conditional expectation $\E[\bm\tau\mid\bm y]$,
which for a multivariate normal is the familiar Bayesian-GLS / ridge
estimator.  Equivalently it is the mode, which solves the
penalized least-squares problem obtained by stacking
\eqref{eq:hp-meas-stacked}--\eqref{eq:hp-state-stacked}:
\begin{equation}
\widehat{\bm\tau}
=\E[\bm\tau\mid\bm y]
=\argmin_{\bm\tau}\ \frac{1}{\sigma_\varepsilon^2}\|\bm y-\bm\tau\|^2
   + \frac{1}{\sigma_e^2}\|D_2\bm\tau\|^2 .
\label{eq:hp-pls-motiv}
\end{equation}
Differentiating and setting the gradient to zero gives a single linear
system,
\begin{equation}
\underbrace{\Bigl(I_T+\lambda\,D_2'D_2\Bigr)}_{J}\,\widehat{\bm\tau}
   = \underbrace{\bm y}_{h},
\qquad \lambda=\frac{\sigma_\varepsilon^2}{\sigma_e^2},
\label{eq:hp-linsys}
\end{equation}
where we divided through by $\sigma_\varepsilon^{-2}$ so that only the
ratio $\lambda$ appears.  This is the closed-form HP solution
$\widehat{\bm\tau}=(I+\lambda D_2'D_2)^{-1}\bm y$ that appears in every
reference --- but now we have \emph{derived} it, as the conditional mean
of a multivariate normal, rather than posited it as an objective.

Three things are worth pausing on, because they are the whole note in
miniature.

\emph{First}, the matrix $J=I+\lambda D_2'D_2$ is \emph{banded}: $D_2$
has bandwidth two, so $D_2'D_2$ is pentadiagonal and $J$ is pentadiagonal.
Solving \eqref{eq:hp-linsys} by a banded Cholesky factorization costs
$O(T)$ operations and $O(T)$ storage, not the $O(T^3)$ a dense
$(I+\lambda D_2'D_2)^{-1}$ would suggest.  The sparsity is not an
accident of the HP filter; it is the stacked image of the fact that
$\tau_t$ interacts only with its near neighbours in time, and it is
exactly the structure that makes the general simultaneous smoother of
this note cheap.

\emph{Second}, the recipe and the model are now visibly the same object.
The HP ``smoothing parameter'' is a variance ratio; the ``smoothing
penalty'' is the $I(2)$ prior; the ``filter'' is a conditional
expectation in a Gaussian.  No forward pass, no backward pass, no
Kalman gain --- just $J\widehat{\bm\tau}=h$, one sparse solve.

\emph{Third}, and most important for what follows, this derivation
generalizes immediately.  We made no use of the specific structure of
the HP filter beyond ``the latent path is constrained by a linear
Gaussian state equation, and the data is a linear Gaussian readout of
the state.''  Any linear Gaussian state-space model has that form, so
any such model reduces to a sparse linear system $J\widehat X=h$ in
exactly the same way.  The rest of this note works out the general
case.

\subsection{What this buys us}

Because the HP filter is now an instance of one Gaussian conditioning
problem rather than a standalone recipe, extensions that would each
require their own derivation in the recursive idiom become small,
composable edits to $J$ and $h$.  A few that we develop later, each
already meaningful for the trend--cycle decomposition above:

\begin{itemize}[leftmargin=2em,itemsep=4pt]

\item \textbf{Exogenous drivers.}  Suppose a credit gap or a terms-of-trade
index $z_t$ shifts trend growth.  Add $\beta z_t$ to the state
equation~\eqref{eq:hp-state} or $\gamma z_t$ to the
measurement~\eqref{eq:hp-meas}.  In simultaneous form this leaves $J$
\emph{unchanged} and only shifts $h$ by a deterministic offset
(Section~\ref{sec:model}).  The smoothed trend then decomposes
additively into a ``pure HP'' part and a credit-driven part.

\item \textbf{Sign and shape constraints.}  Suppose we want to impose
that potential \emph{growth} never turns negative, $g_t=\tau_t-\tau_{t-1}\ge 0$,
or that a NAIRU-like latent process stays non-negative.  In the recursive
form there is no clean way to do this; the Kalman recursion produces an
unconstrained Gaussian.  In simultaneous form it is a single
inequality-constrained quadratic program with the \emph{same} sparse
$J$ as its Hessian, solved by an interior-point method
(Section~\ref{sec:constrained}).

\item \textbf{Estimating $\lambda$ instead of asserting it.}  Since
$\lambda=\sigma_\varepsilon^2/\sigma_e^2$ is a likelihood parameter, we
can estimate it --- and hence let the data choose the degree of
smoothing --- by maximizing the marginal likelihood, which the
simultaneous form computes from the same factorization of $J$ that
delivers the trend (Sections~\ref{sec:loglik} and \ref{sec:em}).

\item \textbf{Robustness to outliers and breaks.}  A one-off oil shock
or a pandemic quarter is a fat-tailed measurement error; a structural
break is a fat-tailed trend innovation.  Replacing the Gaussian errors
with Student-$t$ errors turns \eqref{eq:hp-pls-motiv} into an
iteratively-reweighted version of the same sparse solve
(Section~\ref{sec:robust}), down-weighting the offending observations
automatically.

\item \textbf{Honest uncertainty bands.}  The conditional \emph{variance}
$\Var[\bm\tau\mid\bm y]$ gives pointwise error bands on the trend.  In
the normalization of \eqref{eq:hp-linsys}, where $J=I+\lambda D_2'D_2$
after dividing through by $\sigma_\varepsilon^{-2}$, the posterior
covariance is $\Var[\bm\tau\mid\bm y]=\sigma_\varepsilon^2 J^{-1}$ ---
the measurement-noise scale multiplies the normalized inverse (only if
one imposes $\sigma_\varepsilon^2=1$ is it $J^{-1}$ itself).  The
scale does not affect the smoothed trend, but it does set the width of
the bands.  These widen automatically where data is
missing and narrow where it is dense, and they are computed without
forming the dense inverse, by a sparse selected-inversion recursion
(Section~\ref{sec:selinv}).

\end{itemize}

Every one of these is a small modification of the same object
$J\widehat X=h$.  None requires a new filter, a new backward pass, or a
re-derived boundary condition.  This composability --- many extensions,
one sparse linear system --- is the practical case for the simultaneous
formulation, on top of the formulation being, in our experience, much
easier to understand.  The remainder of the note makes the general
construction precise and then returns, in the worked applications of
Appendix~\ref{sec:apps}, to exactly the kind of trend--cycle problem
this section started with.

\section{A More Intuitive Way to Derive the Kalman Filter}

The Kalman filter and the Rauch--Tung--Striebel (RTS) smoother\footnote{Often
called simply the ``Kalman smoother'' in economics texts.  We retain the
RTS terminology to keep the filter and smoother distinct.} as they appear
in econometrics textbooks descend, almost unaltered, from \citet{Kalman1960}
and from \citet{RTS1965}.  Those papers solved a specific engineering
problem: online estimation of the state of a vehicle from a stream of
noisy sensor readings, on hardware with kilobytes of memory.  Two
constraints shaped the derivation.

First, the problem was \emph{online} --- meaning it had to be solved in
real time.  A new measurement arrived at each instant and the estimate
had to be updated immediately, without revisiting the past.  A forward
recursion of conditional moments is the right object for that setting;
it is what you would invent if you had to.

Second, the hardware could not store the past.  RAM was scarce and
expensive --- the Apollo Guidance Computer, designed around the same
time, had two kilobytes of erasable memory.  Holding the joint
distribution of an entire latent path was inconceivable.  The forward
recursion was not just elegant; it was the only thing that fit.  The
RTS smoother of 1965 added a single backward pass for the same reason:
a second forward sweep would have required storing more than the memory
budget allowed.

Neither constraint binds in modern macroeconometrics.  We are not doing
online state estimation on a missile.  We have a batch of data, on disk,
and we want the smoothed path, the parameters, posterior bands, and the
marginal likelihood jointly, in one shot.  A laptop holds gigabytes; the
joint precision matrix\footnote{Throughout the note the \emph{precision
matrix} of a Gaussian random vector means the inverse of its
variance--covariance matrix.  Econometrics typically works with
covariances directly; the precision form is what makes the sparsity of
state-space models visible and the solves fast, since the inverse of a
dense covariance can be a sparse matrix and vice versa.} for a model
with $n=20$ states and $T=200$ has on the order of $1.6\times 10^5$
nonzeros,\footnote{The precision matrix is block tridiagonal with
$T+1$ diagonal blocks and $T$ off-diagonal blocks, each of size
$n\times n$.  Counting upper-triangular storage of the symmetric
matrix, the structural nonzero count is $(T+1)n^2+Tn^2=(2T+1)n^2$;
for $T=200$ and $n=20$ this gives $401\times 400=160\,400$.}
which is negligible.

Econometric tradition and methodology differ from those used in
engineering.  Control engineers think in terms
of state machines, predictions, and online updates.  Economists think in
terms of joint distributions, regressions, OLS and GLS, and matrix
algebra.  The recursive derivation is in the wrong idiom for us.  It
asks the student to translate econometric questions --- what is the
posterior mean of $x_t$ given all observations? what is the marginal
log-likelihood? what changes when an observation is missing? --- into a
foreign language about filters and predictions, and then translate the
output back.  The simultaneous derivation does not require this.  It
lives in the native econometric idiom: stack the latent path and the
observations into one big multivariate normal vector, and apply the
standard conditional Gaussian formula.  That formula is, structurally,
the GLS estimator; every economist already knows it.  Conditional
moments, regression interpretations, GLS analogies, penalized least
squares, and likelihood are all visible at a glance.

This is the motivation.  The recursive form is not wrong, and remains
the right tool for genuine online filtering.  But for the kind of batch
work that applied macroeconomists actually do, it is needlessly
indirect, and teaching it as the primary derivation propagates that
indirectness to each new cohort of students.

The benefits of the simultaneous approach extend well beyond ease of
understanding.  By placing everything inside one Gaussian conditioning
problem, it produces a unified framework in which extensions of the
basic Gaussian Kalman filter --- heavy-tailed disturbances,
regime-switching dynamics, inequality constraints, missing data,
diffuse priors, time-varying parameters, exogenous inputs, and EM
estimation of unknown parameters --- become small edits to the same
sparse linear system, rather than separate derivations to be
maintained independently.  This is what makes the framework not just
easier to teach, but easier to extend in practice.

\subsection{Where the simultaneous formulation is already standard}
\label{sec:where-used}

Economists' near-universal use of the recursive form is, by a wide
margin, the historical outlier.  In almost every neighbouring
discipline that handles latent Gaussian fields at scale, the
precision-form / simultaneous representation has been the standard
production tool for years, leveraging the weapons-grade sparse linear
algebra infrastructure --- CHOLMOD, MUMPS, PaStiX, SuiteSparse,
Eigen --- that the numerical-analysis community has built over four
decades.  A non-exhaustive map of where this idiom already lives:

\begin{itemize}[leftmargin=2em,itemsep=4pt]

\item \textbf{Spatial statistics and the INLA ecosystem.}  The
\emph{integrated nested Laplace approximation} \citep{RueMartinoChopin2009}
treats spatial, temporal, and spatio-temporal models as Gaussian
Markov random fields with sparse precision matrices and solves them
with sparse Cholesky.  R-INLA has been the dominant Bayesian
inference tool for latent Gaussian models in epidemiology, ecology,
disease mapping, and spatial econometrics for fifteen years.  The
underlying engine is exactly the $JX=h$ solve developed here.
\citet{RueHeld2005} is the canonical reference and motivates the
representation directly.

\item \textbf{Modern Bayesian samplers.}  When Stan
\citep{Carpenter2017} runs Hamiltonian Monte Carlo on a state-space
model with a Gaussian prior over the latent path, the log-density
and gradient computations are simultaneous-form operations on the
sparse precision matrix; the autodiff system never sees a
forward--backward recursion.  Every Stan model with a latent
Gaussian field is implicitly doing simultaneous-form inference at
production scale.

\item \textbf{Large Bayesian VARs and big-data macro.}
\citet{BanburaModugno2014} use an essentially simultaneous EM
representation for mixed-frequency factor models with arbitrary
missing-data patterns, precisely because the recursive form does
not scale gracefully to large panels.  More recent large-BVAR work
\citep{GiannoneLenzaPrimiceri2015} moves in the same direction:
once the number of variables exceeds the few-dozen range, the
precision-form representation is no longer optional.  This is the
macro setting closest to the present note, and the direction of
travel in the literature is clear.

\item \textbf{Geophysical data assimilation.}  Operational weather
prediction at ECMWF and NCEP, ocean reanalysis, and seismic
inversion all rely on variational data assimilation (\emph{4D-Var}),
which is structurally the same object as the simultaneous smoother:
a quadratic penalty on a state trajectory minimized via sparse
linear solves.  The systems handle state dimensions in the millions
on dedicated HPC hardware; none of them runs a recursive Kalman
filter as the production solver.

\item \textbf{Astrostatistics.}  Large-scale Gaussian process
regression with structured covariances --- satellite-data
calibration, gravitational-wave detection pipelines at LIGO,
exoplanet-search likelihood evaluations --- routinely uses
precision-form representations with sparse Cholesky.  The numerical
libraries underneath (CHOLMOD, PaStiX, MUMPS, SuperLU) are the same
ones that any production state-space code would use.

\item \textbf{Image deblurring and structured inverse problems.}  The
classical Tikhonov regularization of an inverse problem,
$\min_X \|y-AX\|^2 + \lambda \|LX\|^2$, is structurally identical
to the simultaneous smoother with $A=\mathcal H$ and $L$ encoding
the state dynamics.  Production image-processing pipelines (medical
imaging, satellite imagery) have solved problems in this form for
decades.

\item \textbf{Computational statistics, broadly.}
\citet{ChanJeliazkov2009} and \citet{McCauslandMillerPelletier2011}
are early macro--statistics references for the precision-form
representation of state-space inference, including efficient
simulation smoothing for posterior draws.  Both papers benchmark
the simultaneous form against the recursive form and find the
former competitive or faster on realistic problem sizes,
contradicting the folk impression that recursion is the
high-performance choice.

\end{itemize}

The common thread is that none of these communities runs a
hand-coded Riccati recursion in user-level code.  They formulate
the latent-Gaussian inference problem as a sparse linear system and
hand it to a production-grade sparse linear algebra library.  The
recursive Kalman filter, viewed from this perspective, is one
particular factorization scheme --- the block-LDL$^{\top}$ forward
sweep developed in Section~\ref{sec:LDL} --- that an
optimizing sparse solver would discover automatically on a
block-tridiagonal matrix, without needing to be programmed by hand.

What this means for econometric practice is concrete: when one
calls \texttt{J $\backslash$ h} in Octave or \texttt{spsolve(J,\,h)}
in SciPy, or invokes any modern sparse-Cholesky routine, one is
standing on roughly two hundred person-years of work by serious
numerical analysts --- much of it specifically tuned for the kind
of structured matrices that arise in state-space problems.  A
hand-coded recursive Kalman loop in \texttt{.m} or \texttt{.py}
gets none of that benefit.

A calibrated statement of the numerical trade-off is in order,
because the recursive-form libraries
(\texttt{statsmodels.tsa.statespace}, \texttt{KFAS}, EViews,
Stata's \texttt{sspace}) have had decades of careful engineering
invested in making the recursive form robust --- square-root
forms, exact diffuse initializations, Joseph forms, UDU and
Bierman--Thornton variants \citep{Bierman1977} --- precisely
because naive Riccati updates can lose symmetry or positive
definiteness through subtraction of covariance matrices.  Sparse
Cholesky on $J$ avoids that particular failure mode by
construction: $J$ is assembled, not propagated, and the
factorization of an SPD matrix never subtracts covariances.  But
the precision form is not a free lunch.  It is a
\emph{normal-equations} representation --- $J$ contains
$Q^{-1}$, $R^{-1}$, and products such as $F'Q^{-1}F$ --- so its
conditioning degrades quadratically where a square-root
(factored) representation degrades linearly, and it requires
$Q$ and $R$ to be invertible, which models with deterministic
state components achieve only through explicit regularization
(Section~\ref{sec:examples} returns to this).  Square-root and
UDU filters were invented for genuinely ill-conditioned online
problems and retain documented advantages there.  The claim of
this note is therefore deliberately scoped: \emph{for well-posed
batch Gaussian smoothing with an SPD precision matrix} --- which
covers the textbook econometric use cases --- sparse Cholesky on
$J$ is a robust, streamlined, and far more teachable coordinate
frame, with the production sparse-linear-algebra stack behind it,
not an unconditional numerical replacement for every Kalman
variant in every regime.

This is the strongest practical case for the simultaneous
formulation.  The underlying mathematics is identical; the
recursive form remains a valid factorization scheme with its own
documented strengths; but for the batch problems this note is
about, the software stack the rest of computational statistics
has converged on assumes the simultaneous representation, and
economists who move to that idiom inherit a numerical
infrastructure that is far larger, more battle-tested, and faster
on real hardware than what any individual research code can hope
to match.

\section{Model}
\label{sec:model}

Consider the linear Gaussian state-space system
\begin{align}
x_t &= F x_{t-1}+ B z_t + w_t,\quad w_t\sim N(0,Q),\quad t=1,\ldots,T,
\label{eq:state-eq}\\
y_t &= Hx_t+ D z_t + v_t,\quad v_t\sim N(0,R),\quad t=1,\ldots,T,
\label{eq:meas-eq}
\end{align}
with initial condition
\begin{equation}
x_0\sim N(a_0,P_0).
\end{equation}
Here $x_t\in\mathbb R^n$ is latent and $y_t\in\mathbb R^m$ is observed.
The vector $z_t\in\mathbb R^k$ is an observed exogenous input --- a
deterministic forcing term that enters the state via the loading
matrix $B$ (size $n\times k$) and the measurement via $D$ (size
$m\times k$).  Setting $B=0$ and $D=0$ recovers the textbook
no-exogenous case.

The two equations describe the same two-layer structure that organizes
all linear state-space modelling.  Equation~\eqref{eq:state-eq} is the
\emph{state} (or \emph{transition}) equation: the unobserved latent
vector $x_t$ evolves autoregressively according to $F$, driven by
exogenous forcing $Bz_t$ and Gaussian innovations $w_t$.
Equation~\eqref{eq:meas-eq} is the \emph{measurement} (or
\emph{observation}) equation: the latent state is observed only through
the linear combination $Hx_t$, contaminated by Gaussian noise $v_t$.
Many familiar macroeconometric objects fit this template:
\begin{itemize}[leftmargin=2em,topsep=2pt,itemsep=1pt]
\item A univariate AR(1) signal observed with noise has $n=m=1$,
  $F=\phi$, $H=1$;
\item The Hodrick--Prescott filter has $n=2$, $m=1$, with $F$
  encoding a trend and a slope component;
\item A small structural VAR has $n=m$ equal to the number of
  observed variables, with $H=I$ and $R$ small;
\item A dynamic factor model has $n\ll m$, with $H$ a tall matrix of
  factor loadings.
\end{itemize}
Section~\ref{sec:examples} works through four cases in full detail ---
an AR(1) signal in noise, the HP filter, an MA(1) via state
augmentation, and an HP filter driven by an observed financial cycle
--- and the worked applications of Appendix~\ref{sec:apps} include a
structural multi-equation system and a dynamic factor model.

Including exogenous inputs from the outset reflects how state-space
models are actually used in macroeconometrics: an output gap is
estimated jointly with the dynamics of a credit cycle, an
unobserved-components model is augmented with policy variables, a
factor model carries observed drivers alongside the latent factors.
The technical apparatus is identical to the no-exogenous case; the
key observation, made precise in Section~\ref{sec:precision-form}, is that
$z_t$ leaves the precision matrix $J$ entirely unchanged and only
shifts the information vector $h$ deterministically.  This is one of
the cleanest payoffs of the simultaneous formulation: the geometry of
the smoothing problem is invariant to deterministic forcing.

The goal of smoothing is to estimate the full path
\begin{equation}
X=(x_0',x_1',\ldots,x_T')'\in\mathbb R^{(T+1)n}
\end{equation}
using all observations $Y=(y_1',\ldots,y_T')'$ and all exogenous inputs
$Z=(z_1',\ldots,z_T')'$.

\subsection{The state dynamics as one block-bidiagonal system}

Define the stacked innovation vector
\begin{equation}
W=\begin{pmatrix}x_0-a_0\\ w_1\\ \vdots\\ w_T\end{pmatrix},
\qquad
\Omega=\Var(W)=\diag(P_0,Q,\ldots,Q).
\end{equation}
Now define
\begin{equation}
G=\begin{pmatrix}
I & 0 & 0 & \cdots & 0\\
-F & I & 0 & \cdots & 0\\
0 & -F & I & \ddots & \vdots\\
\vdots & \ddots & \ddots & \ddots & 0\\
0 & 0 & \cdots & -F & I
\end{pmatrix},
\qquad
c_W=\begin{pmatrix}a_0\\ Bz_1\\ Bz_2\\ \vdots\\ Bz_T\end{pmatrix}.
\end{equation}
The vector $c_W$ collects two things: the prior mean $a_0$ in the
top block, and the deterministic exogenous contributions $Bz_t$ to
each state transition.  With $B=0$ the exogenous blocks vanish and we
recover the no-exogenous case.
Then the full state evolution is simply
\begin{equation}
GX-c_W=W,\qquad W\sim N(0,\Omega).
\label{eq:dyn}
\end{equation}
This is the first key simplification.  The dynamics are not a mysterious
recursion.  They are one sparse linear equation.

\subsubsection{Nilpotent-shift representation}

Let $S_{\downarrow}$ be the lower shift matrix with ones on the first
subdiagonal.  Then
\begin{equation}
G=I_{T+1}\otimes I_n-S_{\downarrow}\otimes F.
\end{equation}
The intuition behind this Kronecker decomposition is that $G$ is the
identity minus a one-step-delay shift weighted by $F$: the diagonal of
$G$ holds the identity on each $x_t$, and the subdiagonal applies
$-F$ to map $x_{t-1}$ to its contribution in the $t$-th equation.
Powers of $S_{\downarrow}$ shift down by more positions, and
$S_{\downarrow}^{T+1}=0$ because one cannot shift further than the
horizon $T+1$.  The Neumann series for $G^{-1}$ therefore terminates
at order $T$:
\begin{equation}
G^{-1}=\sum_{k=0}^{T} S_{\downarrow}^k\otimes F^k.
\end{equation}
The $(t,s)$ block of $G^{-1}$ is therefore
\begin{equation}
(G^{-1})_{ts}=\begin{cases}F^{t-s},&t\ge s,\\0,&t<s.\end{cases}
\end{equation}
This matrix is the explicit propagation operator: shocks are pushed
forward by powers of $F$, and the lower-triangular block structure
encodes the causal arrow of time.

\subsection{The measurement equation}

Stack observations and exogenous inputs as
\begin{equation}
Y=\begin{pmatrix}y_1\\ \vdots\\ y_T\end{pmatrix},
\qquad
Z=\begin{pmatrix}z_1\\ \vdots\\ z_T\end{pmatrix}.
\end{equation}
Define
\begin{equation}
\mathcal H=\begin{pmatrix}
0 & H & 0 & \cdots & 0\\
0 & 0 & H & \cdots & 0\\
\vdots & \vdots & \vdots & \ddots & \vdots\\
0 & 0 & 0 & \cdots & H
\end{pmatrix},\qquad
\mathcal D=\diag(D,D,\ldots,D),\qquad
\Sigma=I_T\otimes R.
\end{equation}
Then
\begin{equation}
Y=\mathcal HX+\mathcal D Z+V,\qquad V\sim N(0,\Sigma).
\label{eq:meas}
\end{equation}
The first block-column of $\mathcal H$ is identically zero: $x_0$ never
enters any observation equation, since measurements run from $t=1$.
The initial state is constrained only by the prior $x_0\sim N(a_0,P_0)$,
which contributes the $P_0^{-1}$ term to the top-left block $J_{00}$ of
the precision matrix.  No $H'R^{-1}H$ term ever lands there.

\subsection{The stacked model as one big multivariate normal}
\label{sec:mvn}

We now have the entire model in two lines:
\begin{align}
GX-c_W &= W,\\
Y - \mathcal HX - \mathcal D Z &= V,
\end{align}
with $W\sim N(0,\Omega)$, $V\sim N(0,\Sigma)$, and $W$ independent of $V$.
Both equations are linear, and the right-hand sides are jointly Gaussian.
The exogenous input $Z$ is observed and deterministic, so it enters only
through means and not variances.  Therefore $(X,Y)$ is itself a
multivariate normal vector of dimension $(T+1)n+Tm$, and its mean and
covariance follow by inspection.

The marginal mean of $X$ is the deterministic propagation of the initial
mean and the exogenous forcing,
\begin{equation}
a=\E X=G^{-1}c_W,
\qquad\text{i.e.}\qquad a_t=F^t a_0+\sum_{s=1}^{t} F^{t-s} B z_s,\ t=0,\ldots,T.
\end{equation}
The first term is the textbook propagation; the second is the
controllability-style integral that maps the exogenous history into the
state-mean path.  The marginal covariance of $X$ is the same as in the
no-exogenous case,
\begin{equation}
\Sigma_X=\Var X=G^{-1}\Omega G^{-\top},
\end{equation}
because $Z$ is observed and deterministic, so it contributes only to
the mean.
The joint mean and covariance of $(X,Y)$ are then
\begin{equation}
\E\begin{pmatrix}X\\Y\end{pmatrix}=\begin{pmatrix}a\\ \mathcal H a+\mathcal D Z\end{pmatrix},
\quad
\Var\begin{pmatrix}X\\Y\end{pmatrix}=
\begin{pmatrix}\Sigma_X & \Sigma_X\mathcal H^\top\\
\mathcal H\Sigma_X & \mathcal H\Sigma_X\mathcal H^\top+\Sigma\end{pmatrix}.
\label{eq:joint}
\end{equation}
The covariance is identical to the no-exogenous case; only the mean
acquires the deterministic offset from $Z$.  That is the whole model.

It is worth pausing to emphasize: equation~\eqref{eq:joint} is the
fundamental object of this entire note.  Everything that follows ---
the Kalman filter, the RTS smoother, the marginal log-likelihood,
EM, the precision-form sparsity, missing data, diffuse priors,
robust extensions, regime-switching, constrained smoothing --- is a
corollary of~\eqref{eq:joint} obtained by applying the standard
conditional-Gaussian formula or by reorganizing the computation to
exploit sparsity.  In particular, the usual econometric intuition
takes over immediately: $X$ is unobserved, $Y$ is observed, and the
operation that recovers the unobserved variable is regression of
$X$ on $Y$.  Carrying out that regression on the joint normal
in~\eqref{eq:joint} is precisely what produces the Kalman smoother
in the next subsection.

\subsubsection{Four examples to ground the abstraction}
\label{sec:examples}

Before applying the conditioning formula, four classical models make
the abstract $G$, $\Omega$, $\mathcal H$, $\Sigma$ concrete.  Each is a
textbook state-space system written first in its conventional
recursive form and then in our stacked matrix form.  The mapping is
mechanical: pick $F,H,Q,R$ from the recursive equations, and the four
big matrices follow.

\paragraph{Example 1: AR(1) signal plus noise.}

The simplest non-trivial state-space model is a scalar AR(1) latent
process observed with noise,
\begin{align}
x_t &= \phi x_{t-1}+w_t,
        \quad w_t\sim N(0,\sigma_w^2),\\
y_t &= x_t+v_t,
        \quad v_t\sim N(0,\sigma_v^2).
\end{align}
All four state-space matrices are scalars: $F=\phi$, $H=1$,
$Q=\sigma_w^2$, $R=\sigma_v^2$.  The dynamics matrix is bidiagonal and
the measurement matrix is a shift of the identity:
\begin{equation}
G=\begin{pmatrix}
1 & 0 & 0 & \cdots & 0\\
-\phi & 1 & 0 & \cdots & 0\\
0 & -\phi & 1 & \ddots & \vdots\\
\vdots & \ddots & \ddots & \ddots & 0\\
0 & 0 & \cdots & -\phi & 1
\end{pmatrix},
\quad
\mathcal H=\begin{pmatrix}
0 & 1 & 0 & \cdots & 0\\
0 & 0 & 1 & \cdots & 0\\
\vdots & \vdots & \vdots & \ddots & \vdots\\
0 & 0 & 0 & \cdots & 1
\end{pmatrix},
\end{equation}
with $\Omega=\diag(P_0,\sigma_w^2,\ldots,\sigma_w^2)$ and
$\Sigma=\sigma_v^2 I_T$.

It is worth working through how the precision matrix arises from the
covariance in this scalar case --- the same logic applies in every
later example, just with $n\times n$ blocks instead of scalars.  From
$X = G^{-1}(W + \text{const})$ the marginal covariance of $X$ is
$\Sigma_X = G^{-1}\Omega G^{-\top}$, which is dense (it is the
classical lag-correlation matrix of an AR(1) plus a top-corner term
from the prior).  Its inverse, however, is sparse:
\begin{equation}
\Sigma_X^{-1} = G^\top\Omega^{-1} G,
\label{eq:precision-from-cov}
\end{equation}
and since $G$ is bidiagonal and $\Omega^{-1}$ is diagonal, the product
$G^\top\Omega^{-1}G$ is tridiagonal.  Adding $\mathcal H^\top\Sigma^{-1}\mathcal H$,
which is itself diagonal because measurements do not couple
periods, gives
\begin{equation}
J = G^\top\Omega^{-1}G + \mathcal H^\top\Sigma^{-1}\mathcal H.
\end{equation}
Carrying out the multiplication block by block yields the explicit
form
\begin{equation}
J=\begin{pmatrix}
P_0^{-1}+\tfrac{\phi^2}{\sigma_w^2} & -\tfrac{\phi}{\sigma_w^2} & & & \\
-\tfrac{\phi}{\sigma_w^2} & \tfrac{1+\phi^2}{\sigma_w^2}+\tfrac{1}{\sigma_v^2} & -\tfrac{\phi}{\sigma_w^2} & & \\
 & \ddots & \ddots & \ddots & \\
 & & -\tfrac{\phi}{\sigma_w^2} & \tfrac{1+\phi^2}{\sigma_w^2}+\tfrac{1}{\sigma_v^2} & -\tfrac{\phi}{\sigma_w^2}\\
 & & & -\tfrac{\phi}{\sigma_w^2} & \tfrac{1}{\sigma_w^2}+\tfrac{1}{\sigma_v^2}
\end{pmatrix},
\end{equation}
with right-hand side $h_0=P_0^{-1}a_0$ and
$h_t=y_t/\sigma_v^2$ for $t\ge 1$.  The smoother is a single
tridiagonal solve, which is what Thomas's algorithm or the Kalman
filter does in this case.  Each diagonal entry has a clear
interpretation: the prior precision plus the precision of the
incoming and outgoing innovations plus the measurement precision.
Smoothing is exactly penalized least squares with state innovations
penalized at scale $1/\sigma_w^2$ and measurement residuals at scale
$1/\sigma_v^2$.  The familiar Kalman smoother formulas drop out of
this tridiagonal solve.

The general lesson: a dense covariance does not mean a dense
calculation, because the inverse of a dense covariance can be a sparse
precision and vice versa.  Equation~\eqref{eq:precision-from-cov} is
the algebraic fact that makes the simultaneous formulation
computationally attractive.

\paragraph{Example 2: the Hodrick--Prescott filter.}

The HP filter has a two-dimensional latent state: the trend level
$\tau_t$ and its slope $\Delta\tau_t$.  The recursive form is
\begin{align}
\tau_t        &= \tau_{t-1}+\Delta\tau_{t-1},\\
\Delta\tau_t  &= \Delta\tau_{t-1}+e_t,\quad e_t\sim N(0,1/\lambda),\\
y_t           &= \tau_t+v_t,\qquad v_t\sim N(0,1).
\end{align}

Before turning to the state-space matrices, it is worth substituting
through the recursion to see what economic content this puts on the
trend.  From the first equation, $\tau_t-\tau_{t-1}=\Delta\tau_{t-1}$;
applying the same identity at $t-1$ gives
$\tau_{t-1}-\tau_{t-2}=\Delta\tau_{t-2}$.  Subtracting,
\begin{equation}
\Delta^2\tau_t \;=\; (\tau_t-\tau_{t-1})-(\tau_{t-1}-\tau_{t-2})
\;=\; \Delta\tau_{t-1}-\Delta\tau_{t-2}
\;=\; e_{t-1} \;\sim\; N(0,\,1/\lambda).
\end{equation}
In other words, the HP filter is the maximum-likelihood smoother for
a model in which the \emph{second difference} of the trend is iid
Gaussian.  When the observed series is interpreted as log output,
this means the HP model assumes that potential log GDP is $I(2)$:
the trend itself is not stationary, its first difference (trend
growth) is not stationary, but the second difference (the change in
trend growth) is white noise.  This is a strong and testable
economic assumption, and it is the right vantage point from which to
ask whether the HP filter is the right model for a given series ---
if the data suggest that potential GDP growth itself is stationary
($I(1)$ rather than $I(2)$ in level), the HP model is misspecified.

With state $x_t=(\tau_t,\Delta\tau_t)^\top$ (dimension $n=2$), the
state-space matrices are
\begin{equation}
F=\begin{pmatrix}1 & 1\\ 0 & 1\end{pmatrix},
\quad
H=\begin{pmatrix}1 & 0\end{pmatrix},
\quad
Q=\begin{pmatrix}\varepsilon & 0\\ 0 & 1/\lambda\end{pmatrix},
\quad
R=1,
\end{equation}
where $\varepsilon>0$ is a tiny regularizer on the level dimension of
$Q$.  This deserves a comment, because it has caused confusion in
practice.  The HP model genuinely has \emph{zero} innovation on the
level component --- $\tau_t$ is determined exactly by $\tau_{t-1}$ and
the slope $\Delta\tau_{t-1}$, with no own shock --- so the natural $Q$
is $\diag(0,1/\lambda)$, singular.  Whether singular $Q$ is a problem
depends on which representation of the smoother one uses.

In the conditioning form, equations~\eqref{eq:smooth-mean}--\eqref{eq:smooth-cov},
the only matrices that need inverting are $P_0$ and $\mathcal H\Sigma_X
\mathcal H^\top+\Sigma$.  Neither requires $Q^{-1}$; the joint
distribution~\eqref{eq:joint} is perfectly well-defined when $Q$ is
singular, since $\Sigma_X=G^{-1}\Omega G^{-\top}$ exists whether or not
$\Omega$ is invertible.  In this representation, $\varepsilon$ is not
needed.

In the \emph{precision} form (Section~\ref{sec:precision-form} below),
we work with $J=G^\top\Omega^{-1}G+\mathcal H^\top\Sigma^{-1}\mathcal H$,
which requires $\Omega^{-1}$ and therefore $Q^{-1}$.  Here the
$\varepsilon>0$ is needed --- not as a modelling assumption, but as a
numerical device to make $Q$ invertible.  As $\varepsilon\to 0$ the
precision-form smoother approaches the singular-$Q$ smoother of the
conditioning form, and in practice any $\varepsilon$ many orders of
magnitude below $1/\lambda$ is invisible in the smoothed path.

Honesty requires spelling out what the regularizer does and does not
preserve.  Inserting $\varepsilon$ changes the \emph{model}: the
precision-form likelihood of Section~\ref{sec:loglik}, the EM bound of
Section~\ref{sec:em}, and the conditioning of $J$ are those of the
regularized model, not of the exact degenerate-Gaussian model with
singular $Q$.  For the smoothed path the difference vanishes as
$\varepsilon\to 0$; for likelihood values it shows up as a
$\log\varepsilon$ term in $\log\det\Omega$ that is \emph{constant} in
the remaining parameters --- so likelihood \emph{ratios}, EM
monotonicity, and ML estimates of the non-degenerate parameters are
unaffected so long as $\varepsilon$ is held fixed, but the likelihood
\emph{level} is that of the regularized model.  Exact alternatives
exist where the distinction matters: eliminate the deterministic state
components analytically before assembling $J$ (for the HP filter this
is precisely the reduction to the trend-only system
$I+\lambda D_2^\top D_2$ of Section~\ref{sec:motivation}), or write
the degenerate-Gaussian likelihood with pseudo-determinants on the
support of $Q$.  The formal statements of this note --- the
LDL$^\top$/Kalman equivalence, the marginal-likelihood formula, the
EM monotonicity --- are claims about models with symmetric
positive-definite $P_0$, $Q$, and $R$; degenerate components are
handled either exactly, by elimination, or approximately, by an
explicit regularizer whose role is disclosed.  The recursive
HP smoother makes the same choice implicitly, hidden inside the
Riccati recursion; the simultaneous form merely makes the
regularization visible because $Q^{-1}$ appears in $J$ literally.

The structure of $G$, $\mathcal H$, $\Omega$, $\Sigma$ is the same as
in Example~1 except that every entry is now a $2\times 2$ block
instead of a scalar.  The precision matrix $J$ is block tridiagonal
with $2\times 2$ blocks.  Concentrating out $\Delta\tau_t$ and writing
the objective in terms of the trend alone, the Gaussian smoothing
problem becomes
\begin{equation}
\min_{\tau}\ \sum_{t=1}^T (y_t-\tau_t)^2
   + \lambda \sum_{t=3}^{T}(\Delta^2\tau_t)^2,
\label{eq:hp-pls}
\end{equation}
which is the original Hodrick--Prescott penalized least squares with
smoothing parameter $\lambda$.  (With the convention
$\Delta^2\tau_t=\tau_t-2\tau_{t-1}+\tau_{t-2}$ adopted above, the
penalty runs over $t=3,\ldots,T$, so that every term involves only
$\tau_1,\ldots,\tau_T$; this is the same $T-2$ curvature terms as the
equivalent indexing $\sum_{t=2}^{T-1}(\tau_{t+1}-2\tau_t+\tau_{t-1})^2$
used in Section~\ref{sec:motivation}, and matches the second-difference
operator $D_2^\top D_2$.)  The classical closed-form solution is
$\widehat\tau=(I+\lambda D_2^\top D_2)^{-1}y$, where $D_2$ is the
second-difference operator.

This is the general lesson: multi-dimensional state is mechanically
no harder than the scalar case.  The blocks of $J$ are $n\times n$
instead of $1\times 1$; the sparsity pattern and the solve are
identical.  And the HP filter, the Beveridge--Nelson decomposition,
unobserved-components models, and the textbook state-space problem
are all the same kind of object --- a quadratic penalty on a latent
path solved by $J\widehat X=h$.  The simultaneous formulation makes
that visible.

\paragraph{Example 3: MA(1) via state augmentation.}

The moving-average model
\begin{equation}
y_t=\epsilon_t+\theta\epsilon_{t-1},
\qquad \epsilon_t\sim N(0,\sigma^2),
\end{equation}
does not fit our framework directly, because $y_t$ depends on both
$\epsilon_t$ and $\epsilon_{t-1}$ while our $\mathcal H$ is block
diagonal in the state.  The textbook fix is to augment the state to
carry the lagged innovation:
\begin{equation}
\label{eq:ma-state}
x_t=\begin{pmatrix}\epsilon_t\\ \epsilon_{t-1}\end{pmatrix},
\quad
F=\begin{pmatrix}0 & 0\\ 1 & 0\end{pmatrix},
\quad
H=\begin{pmatrix}1 & \theta\end{pmatrix}.
\end{equation}
The dynamics inject the new innovation into the first component and
shift the previous one into the second:
$x_t=Fx_{t-1}+(\epsilon_t,0)^\top$, with
$Q=\diag(\sigma^2,\varepsilon)$ for a small regularizer
$\varepsilon$.  Measurement is essentially noiseless,
$R=\varepsilon'$, with a small regularizer for invertibility.

The MA(1) example illustrates a general principle that applies to any
$\mathrm{ARMA}(p,q)$ process: by augmenting the state to dimension
$n=\max(p,q+1)$, the process fits the framework with $F$ in companion
form and $Q$ injecting innovations into the first component.  The
sparsity of $J$ depends on $n$, not on the order of the moving
average.  Long-MA processes get compact sparse formulations.  The
recursive filter applied to this augmented system is exactly what
econometrics textbooks call ``the Kalman filter for ARMA models''.

A subtle point worth flagging now and resolving later: the augmented
state in~\eqref{eq:ma-state} is one of many equivalent realizations of
the same observed MA(1) process.  An MA(1) is identified by the scalar
$\theta$ alone, but our state representation is identified only up to
invertible transformations $x_t\mapsto Mx_t$ with corresponding
adjustments to $F$, $H$, and $Q$.  This rotational ambiguity is
generic; we return to it systematically in Section~\ref{sec:ident}.
For now it suffices to note that any sensible estimation procedure ---
the EM algorithm of Section~\ref{sec:em} in particular --- can only
identify the observable invariants of $(F,H,Q,R)$, not their entries.

\paragraph{Example 4: HP filter with a financial cycle driver.}

The previous examples had no exogenous inputs.  Now consider an
applied case where the latent trend in output is pushed by an
observed financial cycle, say a credit gap $z_t$.  The recursive form
is
\begin{align}
\tau_t        &= \tau_{t-1}+\Delta\tau_{t-1},\\
\Delta\tau_t  &= \Delta\tau_{t-1}+\beta z_t+e_t,\quad e_t\sim N(0,\sigma_e^2),\\
y_t           &= \tau_t+\gamma z_t+v_t,\quad v_t\sim N(0,\sigma_v^2),
\end{align}
where $\beta$ measures how a credit boom raises the trend growth rate
of output, and $\gamma$ is a small direct effect of credit on the
output measurement (a contemporaneous spillover).  The state-space
matrices are
\begin{equation}
F=\begin{pmatrix}1 & 1\\ 0 & 1\end{pmatrix},
\quad
H=\begin{pmatrix}1 & 0\end{pmatrix},
\quad
B=\begin{pmatrix}0\\ \beta\end{pmatrix},
\quad
D=\gamma,
\end{equation}
with $Q=\diag(\varepsilon,\sigma_e^2)$ and $R=\sigma_v^2$.  The
exogenous loading $B$ injects credit-gap shocks into the slope
equation; $D$ allows a direct effect on the observation.

In the stacked formulation, the dynamics matrix $G$ is unchanged from
the HP filter case --- it depends only on $F$.  What changes is the
forcing vector,
\begin{equation}
c_W=\begin{pmatrix}a_0\\ Bz_1\\ \vdots\\ Bz_T\end{pmatrix}=
\begin{pmatrix}a_0\\ (0,\beta z_1)^\top\\ \vdots\\ (0,\beta z_T)^\top\end{pmatrix},
\end{equation}
and the measurement equation acquires the offset $\mathcal D Z=\gamma
(z_1,\ldots,z_T)^\top$.  Conditional on $Z$, the precision matrix $J$
of the smoothing problem is \emph{exactly} the same block-tridiagonal
matrix as in the no-exogenous HP filter; only $h$ shifts to absorb
the deterministic forcing.  The smoothed trend $\widehat\tau_{t\mid T}$
therefore decomposes additively into the no-exogenous HP trend plus
the credit-driven contribution.

This is the practical use case the simultaneous formulation handles
most cleanly.  Coupled with the EM machinery of
Section~\ref{sec:em}, $\beta$ and $\gamma$ can themselves be
estimated jointly with $F,H,Q,R$.  The demo
\texttt{demo\_hp\_credit.m} does exactly that on simulated quarterly
output data with a structural break and outlier contamination, and is
the flagship example of the package.

The four examples make the same point.  Once $F,H,Q,R,B,D$ are
identified from the recursive equations, the abstract framework
becomes a concrete sparse linear-algebra problem.  Different models
just have different sparsity patterns in $J$, and the conditional
Gaussian formula of the next subsection applies to all of them
uniformly.

\subsection{Smoothing as Gaussian conditioning}

Apply to~\eqref{eq:joint} the textbook formula for the conditional
distribution of one block of a multivariate normal given another.
Loosely speaking, what comes out is just a regression --- of the entire
latent path $X$ on all the observations $Y$, with the prior on $X$
playing the role of a regularizer.  The formula is
\begin{equation}
X\mid Y\sim N\bigl(\mu_{X\mid Y},\,\Sigma_{X\mid Y}\bigr)
\end{equation}
with
\begin{align}
\mu_{X\mid Y} &= a+\Sigma_X\mathcal H^\top
   \bigl(\mathcal H\Sigma_X\mathcal H^\top+\Sigma\bigr)^{-1}(Y-\mathcal Ha-\mathcal D Z),
\label{eq:smooth-mean}\\[2pt]
\Sigma_{X\mid Y} &= \Sigma_X-\Sigma_X\mathcal H^\top
   \bigl(\mathcal H\Sigma_X\mathcal H^\top+\Sigma\bigr)^{-1}
   \mathcal H\Sigma_X.
\label{eq:smooth-cov}
\end{align}

That is the Kalman smoother.  There is no recursion, no
prediction/update split, no forward sweep.  It is exactly the Bayesian
GLS regression formula every economist learns in the first semester of
graduate econometrics: read $\mathcal H$ as the regressor matrix, $Y-\mathcal D Z$
as the regressand (the residual after removing the deterministic
effect of $Z$), $\Sigma$ as the error covariance, and $\Sigma_X^{-1}$
as the prior precision on $X$.  The smoothed path $\mu_{X\mid Y}$ is
the posterior mean of this regression; $\Sigma_{X\mid Y}$ is its
posterior covariance.  Smoothing the latent state is the same kind of
object as smoothing the coefficients of a Bayesian regression --- it
just happens to be a regression on a vector that is very long.

The exogenous input $Z$ enters only via the mean of the regressand,
not the covariance, because $Z$ is observed and non-random.

Equations~\eqref{eq:smooth-mean}--\eqref{eq:smooth-cov} contain every
quantity the smoother produces.  Marginal sub-vectors of $\mu_{X\mid Y}$
are the smoothed states $\widehat x_{t\mid T}$.  Diagonal blocks of
$\Sigma_{X\mid Y}$ are the smoothed covariances $P_{t\mid T}$.  Adjacent
off-diagonal blocks are the lag-one covariances $P_{t,t-1\mid T}$ that
EM requires.  The marginal log-likelihood of $Y$ is read off from
the marginal Gaussian on the second block, which has mean $\mathcal Ha+\mathcal D Z$
and covariance $\mathcal H\Sigma_X\mathcal H^\top+\Sigma$.

This is the punch line of the simultaneous formulation: the Kalman
smoother is a special case of the multivariate normal conditioning
formula.  Everything that looks complicated in the recursive derivation
comes from refusing to write the joint distribution down.

\subsubsection{Where each econometric idea sits}

The MVN view turns the rest of this note into a map of standard
econometric machinery onto one Gaussian conditioning problem.  Five
items will reappear later as small edits to either the joint
covariance~\eqref{eq:joint} or the precision matrix $J$ defined in the
next section.
\begin{itemize}[leftmargin=2em]
\item \textbf{GLS.}  Equation~\eqref{eq:smooth-mean} is structurally
  the Bayesian GLS formula: regressor $\mathcal H$, regressand
  $Y-\mathcal Ha$, error covariance
  $\mathcal H\Sigma_X\mathcal H^\top+\Sigma$, prior precision
  $\Sigma_X^{-1}$.  Smoothing \emph{is} Bayesian GLS for the
  regression $Y=\mathcal HX+V$ with prior $X\sim N(a,\Sigma_X)$.
\item \textbf{Marginal log-likelihood.}  The marginal distribution of
  $Y$ in~\eqref{eq:joint} is Gaussian with mean $\mathcal Ha$ and
  covariance $\mathcal H\Sigma_X\mathcal H^\top+\Sigma$, so
  $\log p(Y;\theta)$ is the concentrated GLS criterion: a quadratic
  plus a log-determinant.  Section~\ref{sec:loglik} computes both
  terms without forming that dense covariance.
\item \textbf{Missing data.}  A missing entry of $y_t$ is one row of
  $\mathcal H$ together with the corresponding row and column of
  $\Sigma$ dropped.  The joint MVN is still well-defined; the
  conditional Gaussian still applies.  In precision form
  (Section~\ref{sec:missing}), this is one term subtracted from the
  diagonal block $J_{tt}$.
\item \textbf{Diffuse prior.}  The $P_0^{-1}\to 0$ limit corresponds
  to a flat prior on $x_0$, equivalently an unbounded marginal
  variance in $\Sigma_X$.  $X\mid Y$ remains well-defined whenever
  the observations identify the path.  In precision form
  (Section~\ref{sec:diffuse}), the prior block of $J$ is simply
  dropped.
\item \textbf{EM sufficient statistics.}  The E-step needs
  $\E[x_t x_t'\mid Y]$ and $\E[x_t x_{t-1}'\mid Y]$.  These are
  diagonal and adjacent off-diagonal blocks of
  $\Sigma_{X\mid Y}$ in~\eqref{eq:smooth-cov}, plus outer products
  of smoothed means; Section~\ref{sec:em} reads them off.
\end{itemize}
Each item is the same MVN object viewed through a different
econometric lens.  The next sections introduce the precision form not
because it is conceptually different but because it computes these
same quantities efficiently on a very long Gaussian.

\subsubsection{Why we don't compute estimates directly}

Equations~\eqref{eq:smooth-mean}--\eqref{eq:smooth-cov} are
mathematically perfect, but evaluating them as written --- with dense
matrix algebra --- is computationally wasteful.  Three large dense
objects appear:
\begin{itemize}[leftmargin=2em]
\item $\Sigma_X=G^{-1}\Omega G^{-\top}$ is dense, since $G^{-1}$ is
  lower triangular but full;
\item $\mathcal H\Sigma_X\mathcal H^\top+\Sigma$ is dense of size
  $mT\times mT$;
\item inverting it naively costs $O(m^3T^3)$.
\end{itemize}
This is the cost of treating the smoothing problem as a generic
regression on a vector of length $(T+1)n$.  The next several sections
show that the same calculation can be organized to exploit the
state-space structure: although $\Sigma_{X\mid Y}$ is dense, its
inverse $J=\Sigma_{X\mid Y}^{-1}$ is sparse and block tridiagonal.
Working with $J$ instead of $\Sigma_{X\mid Y}$ reduces the cost from
$O(T^3n^3)$ to $O(Tn^3)$.  No new mathematics is involved --- only a
change of representation that makes the sparsity visible.

\subsection{Identification}
\label{sec:ident}

Identification is not optional.  For any nonsingular $M$,
\begin{equation}
\tilde x_t=Mx_t,
\qquad
\tilde F=MFM^{-1},
\qquad
\tilde H=HM^{-1},
\qquad
\tilde Q=MQM'
\end{equation}
produce the same observed process.  Therefore the latent state is only
identified up to invertible transformations unless restrictions are
imposed.  Common restrictions include:
\begin{itemize}[leftmargin=2em]
\item observed-state normalizations, for example fixing part of $H$ to
  $[I\;0]$;
\item factor-model normalizations, such as diagonal or triangular
  loading matrices;
\item setting the scale of the latent shocks, for example $Q=I$;
\item imposing a companion form, zero restrictions, or stability
  restrictions on $F$;
\item economic long-run restrictions or sign restrictions.
\end{itemize}

When comparing estimated $F,H$ to true values, compare \emph{invariants}
of the latent dynamics (eigenvalues of $F$, achieved log-likelihood, or
the impulse responses of $y_t$) rather than entries of $F$ and $H$.

\section{The simultaneous smoother in precision form}
\label{sec:precision-form}

The remainder of this section develops the precision-form
representation of the smoother in full.  The flow is: posterior
precision and information vector (\S\ref{sec:precision-derivation});
the block-tridiagonal sparsity pattern that makes everything cheap
(\S\ref{sec:tridiag}); the equivalent penalized-least-squares
interpretation (\S\ref{sec:penalized}); the block-LDL$^\top$
factorization, which is the Kalman filter in disguise
(\S\ref{sec:ldl}); the Takahashi backward sweep for posterior
covariances, which is the RTS smoother in disguise
(\S\ref{sec:selinv}); the marginal log-likelihood
(\S\ref{sec:marg-loglik}); and the two standard practical adaptations
that come essentially for free in this representation, missing data
(\S\ref{sec:missing}) and diffuse initialization
(\S\ref{sec:diffuse}).

Each of these is an immediate corollary of the conditional-Gaussian
formula~\eqref{eq:smooth-mean}--\eqref{eq:smooth-cov}, recast in the
representation that exposes the state-space sparsity.

\subsection{Posterior precision and information vector}
\label{sec:precision-derivation}

The posterior kernel of $X\mid Y$ is
\begin{equation}
-\tfrac12(GX-c_W)'\Omega^{-1}(GX-c_W)-\tfrac12(Y-\mathcal HX-\mathcal D Z)'\Sigma^{-1}(Y-\mathcal HX-\mathcal D Z).
\end{equation}
Collecting terms in $X$ gives the posterior precision matrix
\begin{equation}
J=G'\Omega^{-1}G+\mathcal H'\Sigma^{-1}\mathcal H
\label{eq:Jdef}
\end{equation}
and information vector
\begin{equation}
h=G'\Omega^{-1}c_W+\mathcal H'\Sigma^{-1}(Y-\mathcal D Z).
\label{eq:hdef}
\end{equation}
Therefore the smoothed mean and posterior mode solve
\begin{equation}
J\widehat X=h.
\label{eq:simks}
\end{equation}
This is the simultaneous Kalman smoother.

Several features of this formula deserve emphasis.  The exogenous
input $Z$ enters \emph{only} the information vector $h$ and never
the precision matrix $J$.  Concretely, the dynamics contribution
$G'\Omega^{-1}c_W$ adds the propagated terms $Q^{-1}Bz_t$ and
$-F'Q^{-1}Bz_t$ to adjacent blocks of $h$, while the measurement
contribution replaces $Y$ with the residual $Y-\mathcal DZ$ in the
right-hand side.  The block-tridiagonal sparsity, the LDL$^\top$
factorization cost, and the selected-inversion machinery are all
identical to the no-exogenous case.  Switching exogenous inputs on
and off changes only the right-hand side of one sparse linear solve.

Equation~\eqref{eq:Jdef} is just the Woodbury rewriting of the dense
formula~\eqref{eq:smooth-cov}.  To see this, recall that
$\Sigma_X^{-1}=G^\top\Omega^{-1}G$ (since $\Sigma_X=G^{-1}\Omega G^{-\top}$),
and apply the matrix inversion lemma to~\eqref{eq:smooth-cov}:
\begin{equation}
\Sigma_{X\mid Y}^{-1}=\Sigma_X^{-1}+\mathcal H^\top\Sigma^{-1}\mathcal H
=G^\top\Omega^{-1}G+\mathcal H^\top\Sigma^{-1}\mathcal H=J.
\end{equation}
Similarly $\widehat X=\mu_{X\mid Y}=J^{-1}h$ reproduces
\eqref{eq:smooth-mean}.  So $J\widehat X=h$ is not a separate derivation
of the smoother; it is the same Gaussian conditioning written in a
representation that exposes the sparsity.

\subsection{The block-tridiagonal structure}
\label{sec:tridiag}

Since $G$ is block bidiagonal, $G'\Omega^{-1}G$ is block tridiagonal.
Since measurements enter one period at a time, $\mathcal H'\Sigma^{-1}\mathcal H$
is block diagonal.  Thus $J$ is block tridiagonal.

For the time-invariant model,
\begin{equation}
\resizebox{\textwidth}{!}{$
\renewcommand{\arraystretch}{1.25}
\setlength{\arraycolsep}{3pt}
J=\begin{pmatrix}
P_0^{-1}+F'Q^{-1}F & -F'Q^{-1} & 0 & \cdots & 0\\
-Q^{-1}F & Q^{-1}+F'Q^{-1}F+H'R^{-1}H & -F'Q^{-1} & \cdots & 0\\
0 & -Q^{-1}F & Q^{-1}+F'Q^{-1}F+H'R^{-1}H & \ddots & \vdots\\
\vdots & \vdots & \ddots & \ddots & -F'Q^{-1}\\
0 & 0 & \cdots & -Q^{-1}F & Q^{-1}+H'R^{-1}H
\end{pmatrix}.
$}
\end{equation}
The corresponding right-hand side is, in the no-exogenous case,
\begin{equation}
h=\begin{pmatrix}P_0^{-1}a_0\\ H'R^{-1}y_1\\ H'R^{-1}y_2\\ \vdots\\ H'R^{-1}y_T\end{pmatrix}.
\end{equation}
With exogenous inputs, the right-hand side becomes
\begin{equation}
h=\begin{pmatrix}P_0^{-1}a_0-F'Q^{-1}Bz_1\\[2pt]
H'R^{-1}(y_1-Dz_1)+Q^{-1}Bz_1-F'Q^{-1}Bz_2\\[2pt]
H'R^{-1}(y_2-Dz_2)+Q^{-1}Bz_2-F'Q^{-1}Bz_3\\[2pt]
\vdots\\[2pt]
H'R^{-1}(y_T-Dz_T)+Q^{-1}Bz_T\end{pmatrix},
\end{equation}
where the $Bz_t$ contributions enter from both the predicting and the
predicted side of each transition.  The block-tridiagonal structure of
$J$ is unaffected.

\subsection{Equivalent penalized least-squares problem}
\label{sec:penalized}

The same object is obtained as the minimizer of
\begin{align}
L(X)
&=(x_0-a_0)'P_0^{-1}(x_0-a_0)\\
&\quad+\sum_{t=1}^T (x_t-Fx_{t-1}-Bz_t)'Q^{-1}(x_t-Fx_{t-1}-Bz_t)\\
&\quad+\sum_{t=1}^T (y_t-Hx_t-Dz_t)'R^{-1}(y_t-Hx_t-Dz_t).
\end{align}
The first-order conditions are exactly $J\widehat X=h$.  This view is
usually the most transparent one: smoothing is a structured penalized
least-squares problem.  The exogenous terms $Bz_t$ and $Dz_t$ enter the
objective as deterministic shifts of the residuals, so they fall out
when the gradient is taken; they appear only on the right-hand side
of the resulting linear system.

\subsection{The recursive filter as a block-LDL$^{\top}$ forward sweep}
\label{sec:ldl}
\label{sec:LDL}

The recursive Kalman filter and the simultaneous smoother are the same
object.  This section makes that statement precise.

Because $J$ is symmetric positive definite and block tridiagonal, it
admits a unique block-LDL$^\top$ factorization
\begin{equation}
J=LDL^\top,
\end{equation}
where $L$ is block lower-bidiagonal with identity diagonal blocks, and
$D=\diag(D_0,D_1,\ldots,D_T)$.  Write the off-diagonal blocks of $L$ as
$L_{t,t-1}=M_t$; this is the convention used by the package code
(\texttt{simks\_selected\_inv}).\footnote{Some references instead set
$L_{t,t-1}=-M_t$, which flips the sign of every $M_t$ below and of the
off-diagonal line in the Takahashi recursion of
Section~\ref{sec:selinv}.  Either convention is fine; mixing them is
not.}

Matching blocks of $LDL^\top$ to blocks of $J$ determines the factors.
The $(t,t-1)$ block of $LDL^\top$ is $L_{t,t-1}D_{t-1}=M_tD_{t-1}$, and
the $(t,t)$ block is $D_t+M_tD_{t-1}M_t^\top$.  Because the blocks do
not commute, the order of multiplication matters: $M_t$ must be
recovered by a \emph{right} solve against $D_{t-1}$,
\begin{align}
D_0 &= J_{00}=P_0^{-1}+F'Q^{-1}F,\\
M_t &= J_{t,t-1}\,D_{t-1}^{-1}=-\,Q^{-1}F\,D_{t-1}^{-1},
\qquad t\ge 1,\\
D_t &= J_{tt}-M_t\,D_{t-1}\,M_t^\top.
\label{eq:LDLstep}
\end{align}

This recursion is the Kalman filter, written in precision form --- but
the identification of the pivots requires one step of care.  The
diagonal pivots $D_t$ are \emph{not} the filtered precisions for
interior periods, because the block $J_{tt}=Q^{-1}+F'Q^{-1}F+H'R^{-1}H$
natively carries the prior-precision contribution $F'Q^{-1}F$ of the
\emph{outgoing} transition to $x_{t+1}$.  Define instead
\begin{equation}
C_t \;=\; D_t-F'Q^{-1}F \quad (t<T),
\qquad
C_T \;=\; D_T,
\label{eq:Ct-def}
\end{equation}
stripping the outgoing term (the terminal block has no outgoing
transition, which is why $D_T$ needs no correction).  Substituting
\eqref{eq:LDLstep} and $J_{tt}$ into \eqref{eq:Ct-def},
\begin{equation}
C_t
= Q^{-1}-Q^{-1}F\,\bigl(C_{t-1}+F'Q^{-1}F\bigr)^{-1}F'Q^{-1}
  + H'R^{-1}H,
\end{equation}
and the Woodbury identity collapses the first two terms:
\begin{equation}
C_t=\bigl(Q+F\,C_{t-1}^{-1}F'\bigr)^{-1}+H'R^{-1}H .
\end{equation}
Setting $P_{t\mid t}=C_t^{-1}$, this is exactly
\begin{equation}
P_{t\mid t}^{-1}=P_{t\mid t-1}^{-1}+H'R^{-1}H,
\qquad
P_{t\mid t-1}=FP_{t-1\mid t-1}F'+Q,
\end{equation}
the filtered/predicted covariance updates of the textbook Kalman
recursion, in their precision-form equivalents, with the anchor
$C_0=D_0-F'Q^{-1}F=P_0^{-1}=P_{0\mid 0}^{-1}$.  The forward solve
$Lu=h$ likewise produces the filtered \emph{means} in information
form: $u_t=C_t\,\widehat x_{t\mid t}=P_{t\mid t}^{-1}\widehat x_{t\mid t}$,
the information state of the information-filter literature.

The backward solve $L^\top \widehat X=D^{-1}u$ then propagates information
backward.  Identifying the smoothing gain
\begin{equation}
A_t=P_{t\mid t}F' P_{t+1\mid t}^{-1}
\end{equation}
in the back-substitution yields
\begin{equation}
\widehat x_{t\mid T}=\widehat x_{t\mid t}+A_t\bigl(\widehat x_{t+1\mid T}-F\widehat x_{t\mid t}\bigr),
\end{equation}
which is the Rauch--Tung--Striebel smoother.

The message is not that ``$J\widehat X=h$ is faster than the recursive
filter''.  Both cost $O(Tn^3)$, because they are the same operation
written two ways.  The message is that the Kalman filter is the
block-LDL$^\top$ factorization of $J$ done by hand, and that this
factorization is what makes the operation $O(Tn^3)$ rather than
$O(T^3n^3)$.  Modern sparse direct solvers (CHOLMOD, etc.) discover
this factorization automatically.  See \citet{Davis2006}.

\subsection{Posterior covariances via selected inversion}
\label{sec:selinv}

The posterior covariance of the smoothed path is $J^{-1}$.  Forming
$J^{-1}$ densely costs $O(T^3n^3)$ and is wasteful: we usually want only
the diagonal and first sub-diagonal blocks
\begin{equation}
\Sigma_{tt}=\Cov(x_t\mid Y),
\qquad
\Sigma_{t+1,t}=\Cov(x_{t+1},x_t\mid Y),
\end{equation}
for EM updates and for posterior bands.

These blocks can be computed by Takahashi's recursion, due to
\citet{TakahashiFaganChen1973} and \citet{ErismanTinney1975}.  Starting from $LDL^\top$
of the previous section, $J^{-1}=L^{-\top}D^{-1}L^{-1}$.  With
$M_t=L_{t,t-1}=J_{t,t-1}D_{t-1}^{-1}$ exactly as defined in
Section~\ref{sec:LDL}, the recursion is, for $t=T,T-1,\ldots,0$,
\begin{align}
\Sigma_{TT}&=D_T^{-1},\\
\Sigma_{t+1,t}&=-\Sigma_{t+1,t+1}M_{t+1},\\
\Sigma_{tt}&=D_t^{-1}+M_{t+1}^\top\,\Sigma_{t+1,t+1}\,M_{t+1}.
\end{align}
(These three lines are derived by equating the strict upper triangle
of $L^\top\Sigma=D^{-1}L^{-1}$ to its right-hand side and reading off
blocks; the sign on the middle line is tied to the convention
$L_{t,t-1}=+M_t$ --- under the alternative convention
$L_{t,t-1}=-M_t$ it flips, a classic source of transcription errors.
Note that for the state-space blocks $M_{t+1}=-Q^{-1}FD_t^{-1}$, so
$\Sigma_{t+1,t}=\Sigma_{t+1,t+1}Q^{-1}FD_t^{-1}$ is positive in the
scalar AR(1) case with $\phi>0$, as it must be: adjacent smoothed
states are positively correlated under persistent dynamics.)
That is: at each step backward, we need the diagonal block just above
and only the LDL factors of $J$.  Total cost $O(Tn^3)$.

The familiar RTS smoother covariance update
\begin{equation}
P_{t\mid T}=P_{t\mid t}+A_t(P_{t+1\mid T}-P_{t+1\mid t})A_t'
\end{equation}
is exactly this Takahashi backward recursion in the precision-form
coordinates of Section~\ref{sec:LDL}.

\subsection{Marginal log-likelihood}
\label{sec:marg-loglik}
\label{sec:loglik}

For maximum likelihood, model comparison, or any kind of likelihood-based
inference, we need $\log p(Y;\theta)$, not the joint $\log p(X,Y;\theta)$
at a particular $X$.  The simultaneous formulation gives it cleanly.

From~\eqref{eq:dyn} and~\eqref{eq:meas}, $(X,Y)$ is jointly Gaussian.
Standard manipulations give
\begin{equation}
\log p(Y;\theta) = -\tfrac12(Y-\mathcal H a-\mathcal D Z)'\Omega_Y^{-1}(Y-\mathcal H a-\mathcal D Z)
                   -\tfrac12 \log\det \Omega_Y
                   -\tfrac{mT}{2}\log(2\pi),
\end{equation}
where $a=G^{-1}c_W$ (which now incorporates the propagated $Bz_t$) and
$\Omega_Y=\mathcal H G^{-1}\Omega G^{-\top}\mathcal H'+\Sigma$.  Working
directly with $\Omega_Y$ is what the recursive Kalman filter does
implicitly via the innovations representation, at $O(Tn^3)$ cost.

The simultaneous form is more direct.  Complete the square in $X$:
\begin{equation}
\log p(X,Y;\theta)
=-\tfrac12 X'JX+h'X+\text{const}(\theta,Y).
\end{equation}
Integrating $X$ out using
$\int\exp(-\tfrac12X'JX+h'X)\,dX
=(2\pi)^{(T+1)n/2}(\det J)^{-1/2}\exp(\tfrac12h'J^{-1}h)$ gives
\begin{equation}
\log p(Y;\theta)
=+\tfrac12 h'J^{-1}h-\tfrac12 \log\det J + \text{const}'(\theta,Y),
\label{eq:marg-ll-integrated}
\end{equation}
where the remaining constant collects the prior and measurement
normalizing terms.  The signs deserve a second look, because both are
easy to get backwards: the completed-square quadratic enters with a
\emph{plus} ($+\tfrac12h'J^{-1}h$), and the log-determinant enters with
a \emph{minus} ($-\tfrac12\log\det J$).  The sanity check on the
latter: a more concentrated posterior on the latent path means a larger
$\det J$, and concentrating the latent states strictly \emph{reduces}
the integrated mass, so $\log\det J$ must lower the marginal
likelihood.  Writing out the constant, and using
$\widehat X=J^{-1}h$ so that $h'J^{-1}h=\widehat X'h$,
\begin{align}
\log p(Y;\theta)
&=\tfrac12 \widehat X'h
 -\tfrac12 c_W'\Omega^{-1}c_W
 -\tfrac12 (Y-\mathcal DZ)'\Sigma^{-1}(Y-\mathcal DZ) \nonumber\\
&\quad -\tfrac12 \log\det J
 -\tfrac12 \log\det\Omega
 -\tfrac12 \log\det\Sigma
 -\tfrac{mT}{2}\log(2\pi).
\label{eq:marg-ll}
\end{align}
Two exogenous-input details matter here and are easy to drop.  First,
the measurement quadratic uses the exogenous-adjusted residual
$Y-\mathcal DZ$, not $Y$ itself.  Second, the dynamics quadratic uses
the full stacked offset $c_W=(a_0',\,(Bz_1)',\ldots,(Bz_T)')'$, so
\begin{equation}
c_W'\Omega^{-1}c_W
=a_0'P_0^{-1}a_0+\sum_{t=1}^T (Bz_t)'Q^{-1}(Bz_t),
\end{equation}
and $h$ in the $\widehat X'h$ term is the full information vector of
\eqref{eq:hdef}, which already carries the propagated $Bz_t$
contributions.  (With $B=D=0$ both reduce to the familiar
no-exogenous expressions $a_0'P_0^{-1}a_0$ and $Y'\Sigma^{-1}Y$.)
Every term in~\eqref{eq:marg-ll} is cheap.  The log-determinant
$\log\det J$ is read off from the diagonal of $D$ in $J=LDL^\top$,
or equivalently from twice the sum of $\log$ of the diagonal of the
Cholesky factor.  Sparse Cholesky on $J$ delivers both $\widehat X$ and
$\log\det J$ in a single call.

\subsection{Missing data}
\label{sec:missing}

Suppose at time $t$ only a subset $S_t\subseteq\{1,\ldots,m\}$ of the
observation channels is available, and $y_{t}^{S_t}$ denotes the
observed sub-vector.  In the simultaneous formulation, the only thing
that changes is the measurement contribution at $t$:
\begin{equation}
H'R^{-1}H \;\longrightarrow\; H_{S_t}'R_{S_t}^{-1}H_{S_t},
\qquad
H'R^{-1}y_t \;\longrightarrow\; H_{S_t}'R_{S_t}^{-1}y_t^{S_t}.
\end{equation}
A fully missing observation contributes nothing; the corresponding
diagonal block of $J$ is just $Q^{-1}+F'Q^{-1}F$ (or its endpoint
analogue), with no measurement term.

There is no special ``missing-data Kalman filter''.  There is just $J$
and $h$, assembled from whatever measurements exist.

\subsection{Diffuse initialization}
\label{sec:diffuse}

If we do not want to commit to a prior on $x_0$, take the limit
$P_0^{-1}\to 0$.  In~\eqref{eq:simks}, this simply drops the prior block
from the top-left of $J$ and the term $P_0^{-1}a_0$ from $h$.  Provided
the dynamics and measurements together identify $x_0$ (essentially: $x_0$
appears in some path that has observations downstream), $J$ remains
positive definite and the solve is well-posed.  This is the
``proper'' diffuse limit; the recursive analogue requires the careful
exact-initial-Kalman-filter treatment of \citet{deJong1989} and
\citet{KoopmanDurbin2003}.

\section{Time-varying systems}
\label{sec:tv}

A third small edit, also forecast in Section~\ref{sec:mvn}.  Replacing
the constant matrices $F,H,Q,R$ with sequences $F_t,H_t,Q_t,R_t$ leaves
the conditional Gaussian~\eqref{eq:smooth-mean}--\eqref{eq:smooth-cov}
formally unchanged; only the entries of $G$, $\Omega$, $\mathcal H$,
and $\Sigma$ change.  For
\begin{align}
x_t&=F_tx_{t-1}+w_t,\quad w_t\sim N(0,Q_t),\\
y_t&=H_tx_t+v_t,\quad v_t\sim N(0,R_t),
\end{align}
the dynamics matrix becomes
\begin{equation}
G=\begin{pmatrix}
I & 0 & 0 & \cdots & 0\\
-F_1 & I & 0 & \cdots & 0\\
0 & -F_2 & I & \ddots & \vdots\\
\vdots & \ddots & \ddots & \ddots & 0\\
0 & 0 & \cdots & -F_T & I
\end{pmatrix}
\end{equation}
and $\Omega=\diag(P_0,Q_1,\ldots,Q_T)$.  The precision matrix remains
block tridiagonal, with blocks
\begin{align}
J_{00} &= P_0^{-1}+F_1'Q_1^{-1}F_1,\\
J_{tt} &= Q_t^{-1}+F_{t+1}'Q_{t+1}^{-1}F_{t+1}+H_t'R_t^{-1}H_t,\quad 1\le t\le T-1,\\
J_{TT} &= Q_T^{-1}+H_T'R_T^{-1}H_T,\\
J_{t,t-1}&=-Q_t^{-1}F_t,\qquad J_{t-1,t}=-F_t'Q_t^{-1}.
\end{align}
Everything downstream --- LDL$^\top$ factorization, selected inversion,
marginal log-likelihood, EM --- carries over unchanged.  This is
the modularity claim in concrete form.

\section{Estimating Model Parameters}
\label{sec:em}

Suppose $F$ and $H$ are unknown but constant over time.  The natural
objective is
\begin{align}
\min_{X,F,H,Q,R} \quad
&(x_0-a_0)'P_0^{-1}(x_0-a_0)\\
&+\sum_{t=1}^T (x_t-Fx_{t-1})'Q^{-1}(x_t-Fx_{t-1})\\
&+\sum_{t=1}^T (y_t-Hx_t)'R^{-1}(y_t-Hx_t),
\end{align}
together with the implicit $-T\log\det Q-T\log\det R$ from the Gaussian
normalizing constants if $Q,R$ are also to be estimated.

\subsection{Quick-and-dirty: alternating regressions}

The simplest scheme is: given $(F,H,Q,R)$, solve for $X$ with the
simultaneous smoother; given $X$, update $F$ and $H$ by ordinary least
squares,
\begin{equation}
\widehat F=\bigl(\textstyle\sum_t x_tx_{t-1}'\bigr)\bigl(\textstyle\sum_t x_{t-1}x_{t-1}'\bigr)^{-1},
\quad
\widehat H=\bigl(\textstyle\sum_t y_tx_t'\bigr)\bigl(\textstyle\sum_t x_tx_t'\bigr)^{-1},
\end{equation}
and then update $Q,R$ by residual covariances.  This is conceptually
clean but the $Q$ and $R$ estimates are biased: they ignore the fact
that the smoothed $x_t$ is a posterior mean, not a sample from the
posterior.

\subsection{Proper Gaussian EM}

The unbiased fix is the standard EM algorithm of \citet{ShumwayStoffer1982}
and \citet{Watson1983}.  The E-step computes the smoothed sufficient
statistics
\begin{align}
S_{00} &= \sum_{t=1}^T \E[x_{t-1}x_{t-1}'\mid Y] = \sum_{t=1}^T \bigl(\widehat x_{t-1\mid T}\widehat x_{t-1\mid T}'+\Sigma_{t-1,t-1}\bigr),\\
S_{11} &= \sum_{t=1}^T \E[x_tx_t'\mid Y] = \sum_{t=1}^T \bigl(\widehat x_{t\mid T}\widehat x_{t\mid T}'+\Sigma_{tt}\bigr),\\
S_{10} &= \sum_{t=1}^T \E[x_tx_{t-1}'\mid Y] = \sum_{t=1}^T \bigl(\widehat x_{t\mid T}\widehat x_{t-1\mid T}'+\Sigma_{t,t-1}\bigr).
\end{align}
The posterior covariance terms $\Sigma_{tt}$ and $\Sigma_{t,t-1}$ are
delivered by the selected-inversion recursion of Section~\ref{sec:selinv}.
The M-step then solves
\begin{align}
\widehat F &= S_{10}S_{00}^{-1},\\
\widehat Q &= \tfrac1T\bigl(S_{11}-S_{10}S_{00}^{-1}S_{10}'\bigr),\\
\widehat H &= \bigl(\textstyle\sum_t y_t\widehat x_{t\mid T}'\bigr)\bigl(\textstyle\sum_t (\widehat x_{t\mid T}\widehat x_{t\mid T}'+\Sigma_{tt})\bigr)^{-1},\\
\widehat R &= \tfrac1T\textstyle\sum_t \bigl((y_t-\widehat H\widehat x_{t\mid T})(y_t-\widehat H\widehat x_{t\mid T})' + \widehat H\Sigma_{tt}\widehat H'\bigr).
\end{align}
The initial-state parameters are updated as well, by the standard
\citet{ShumwayStoffer1982} single-observation M-step
\begin{equation}
\widehat a_0=\widehat x_{0\mid T},
\qquad
\widehat P_0=\Sigma_{00}
\;\bigl(=\Sigma_{00}+(\widehat x_{0\mid T}-\widehat a_0)(\widehat x_{0\mid T}-\widehat a_0)'
\text{ evaluated at the new } \widehat a_0\bigr),
\end{equation}
which is what \texttt{simks\_em\_const} and \texttt{em\_const}
implement; the returned \texttt{a0} and \texttt{P0} are therefore
genuinely estimated quantities, not pass-throughs of the
initialization.  (Estimating $(a_0,P_0)$ from one effective
observation is, of course, weakly informative; fixing them, or using
the diffuse initialization of Section~\ref{sec:diffuse}, are the
standard alternatives.)
The marginal log-likelihood~\eqref{eq:marg-ll} is non-decreasing across
iterations.  Monitoring it is a useful sanity check; the demo
\texttt{demo\_const\_em.m} does exactly that.

The cost per iteration is dominated by the smoother solve and the
selected inversion, both $O(Tn^3)$.

\subsection{Exogenous loadings as a regression augmentation}

When the model includes exogenous inputs, the EM M-step generalizes
in the most natural way possible.  The state-side update becomes a
joint regression of $\widehat x_{t\mid T}$ on the augmented regressor
$(\widehat x_{t-1\mid T},\,z_t)$:
\begin{equation}
\begin{pmatrix}\widehat F & \widehat B\end{pmatrix}
=\begin{pmatrix}S_{10} & S_{1z}\end{pmatrix}
\begin{pmatrix}S_{00} & S_{0z}\\ S_{0z}^\top & S_{zz}\end{pmatrix}^{-1},
\label{eq:em-FB}
\end{equation}
where $S_{1z}=\sum_t \widehat x_{t\mid T}z_t^\top$,
$S_{0z}=\sum_t \widehat x_{t-1\mid T}z_t^\top$, and
$S_{zz}=\sum_t z_t z_t^\top$.  The same construction applies to the
measurement side:
\begin{equation}
\begin{pmatrix}\widehat H & \widehat D\end{pmatrix}
=\begin{pmatrix}S_{yx} & S_{yz}\end{pmatrix}
\begin{pmatrix}S_{xx} & S_{xz}\\ S_{xz}^\top & S_{zz}^{\text{obs}}\end{pmatrix}^{-1},
\end{equation}
and the noise covariances $\widehat Q,\widehat R$ then follow by
the standard residual-sum-of-squares formulas with cross-terms.

The shape of these updates is the substantive payoff of including
$z_t$ from the outset.  Estimating the loading of a credit gap on
output trend is mechanically identical to estimating the autoregressive
coefficient: both are coefficients on a regressor in
$\widehat x_t = [\widehat F\,\widehat B](\widehat x_{t-1},\,z_t)^\top$.
The simultaneous smoother delivers the smoothed sufficient statistics
the regression needs.

The demo \texttt{demo\_hp\_credit.m} runs this machinery on a
simulated HP-plus-credit-cycle setup with structural breaks and
outlier contamination.  To be precise about what is combined and how:
the demo compares \emph{separate} estimators on the same data ---
Gaussian smoothing without and with the true exogenous loadings,
robust Student-$t$ smoothing (Section~\ref{sec:robust}) at the true
parameters, and Gaussian EM that estimates $(F,H,Q,R,B,D)$ jointly.
It does \emph{not} run an integrated robust-EM procedure (alternating
IRLS weight updates with M-step parameter updates); such a generalized
EM is a natural and straightforward extension of the two building
blocks, but it is not implemented in the package, and we flag that
explicitly rather than imply otherwise.  The realistic ingredients for
applied macroeconomic work are nonetheless all present: parameters are
unknown, the data are imperfect, and the underlying dynamics include
observed drivers.

\section{Beyond the Standard Model}
\label{sec:beyond}

So far the smoothing problem has been a textbook linear Gaussian
state-space system with constant or time-varying matrices.  Most
substantively interesting macroeconomic state-space problems
violate at least one of these assumptions: data contain outliers
or structural breaks (so the Gaussian assumption fails); regimes
shift between expansions and recessions (so the linearity
assumption fails); or known structural bounds rule out parts of
the latent path (so the posterior is truncated).  The recursive
Kalman filter loses its closed-form structure in each of these
cases and is patched only by case-specific machinery.  The
simultaneous form, by contrast, handles all three with the same
sparse linear-algebra primitives: in every case the algorithm is
either iteratively reweighted solves on the same $J$, or
variational EM with $J$ assembled from responsibility-weighted
expected quadratic forms, or a sparse QP whose Hessian is $J$ plus
a low-rank update.  This section develops the three cases.

\subsection{Non-Gaussian disturbances}
\label{sec:robust}

The stacked formulation extends naturally beyond Gaussian errors.
Replace the quadratic penalties by general negative log densities:
\begin{equation}
\min_X \sum_{t=1}^T \rho_w(x_t-Fx_{t-1})+
\sum_{t=1}^T \rho_v(y_t-Hx_t)+\rho_0(x_0-a_0).
\end{equation}
The exact closed-form solve is lost, but the sparsity is preserved:
each $x_t$ still interacts only with $x_{t-1}$, $x_{t+1}$, and $y_t$.
The same $J$ pattern, with weights updated each iteration, supports a
broad family of robust state-space smoothers.  This section makes that
concrete.

\subsubsection{Heavy tails as Gaussian scale mixtures}

The practically important case is heavy-tailed measurement noise (to
robustify against outliers) and heavy-tailed state innovations (to
permit structural breaks).  Both are handled cleanly by the
\emph{Gaussian scale mixture} representation
\citep{AndrewsMallows1974}:
\begin{equation}
v_t \mid \tau_t \sim N(0,\,R/\tau_t),
\qquad \tau_t \sim p(\tau \mid \text{family}).
\end{equation}
Marginalizing $\tau_t$ out gives the desired heavy-tailed marginal:
\begin{itemize}[leftmargin=2em]
\item Student-$t$ with $\nu$ degrees of freedom: $\tau \sim
  \text{Gamma}(\nu/2,\,\nu/2)$.
\item Laplace (double exponential): $\tau$ has an inverse-exponential
  distribution.
\end{itemize}
Throughout this section, $\delta_t=(r_t^\top R^{-1}r_t)^{1/2}$ denotes
the \emph{Mahalanobis distance} of the period-$t$ residual, so
$\delta_t^2$ is the squared distance (the quadratic form itself); the
two are easy to confuse in implementation, and the update formulas
below are stated in terms of $\delta_t^2$ where the quadratic form is
meant.  The Huber loss is deliberately \emph{not} on the list above:
its working weight $\min(1,\,c/\delta_t)$ depends on the realized
residual, whereas a mixing distribution $p(\tau)$ in a proper Gaussian
scale mixture cannot.  Huber smoothing is therefore best understood
not as a latent-scale model but as classical IRLS on the Huber
$\rho$-function \citep{HollandWelsch1977}: the weight
$\psi(\delta)/\delta=\min(1,c/\delta)$ is the majorize--minimize
multiplier that makes each reweighted Gaussian solve a tangent
quadratic to the Huber objective.  It plugs into exactly the same
sparse machinery, but its justification is the loss function, not a
mixture representation.
Conditional on the $\tau_t$, the model is Gaussian with time-varying
covariance $R_t = R/\tau_t$, and the simultaneous smoother of
Section~\ref{sec:tv} applies without modification.

\subsubsection{EM and IRLS on the same sparse system}

The natural inference algorithm alternates reweighted Gaussian solves
with weight updates \citep{LangeLittleTaylor1989}.
\begin{enumerate}[leftmargin=2em]
\item \textbf{Reweighted solve.}  Given the current weights $\tau_t^{(k)}$,
  solve the time-varying Gaussian smoother with $R_t = R/\tau_t^{(k)}$.
  This is one sparse Cholesky on the same $J$ pattern.
\item \textbf{Weight update.}  Given the smoothed $\widehat X^{(k+1)}$,
  update each weight from the smoothed residual
  \begin{equation}
  r_t = y_t - H \widehat x_{t\mid T}^{(k+1)} - Dz_t,
  \end{equation}
  where the exogenous offset $Dz_t$ must be subtracted whenever the
  model has measurement-side inputs.  For Student-$t$,
  \begin{equation}
  \tau_t^{(k+1)} = \frac{\nu + m_t}{\nu + r_t^\top R^{-1} r_t},
  \end{equation}
  where $m_t$ is the number of observed entries of $y_t$.  For Laplace,
  $\tau_t = (r_t^\top R^{-1} r_t)^{-1/2}$.  For Huber,
  $\tau_t = \min(1,\, c/\delta_t)$, with the canonical
  $c=1.345$ \citep{HollandWelsch1977}.
\end{enumerate}
Each iteration costs one $O(Tn^3)$ sparse Cholesky.  The sparsity
pattern of $J$ is identical to the Gaussian case; only the numerical
values of the diagonal blocks change.

It is worth being precise about what this iteration computes, because
two related but distinct targets are in circulation.  As written ---
with weights evaluated by \emph{plugging in} the residual at the
current smoothed mean --- the fixed point is the posterior
\emph{mode} of the heavy-tailed model: plug-in IRLS is exactly
majorize--minimize on the non-Gaussian MAP objective.  This is a
perfectly respectable estimand, and it is the package default.  It is
\emph{not}, however, the E-step of an exact latent-scale EM algorithm,
which requires the posterior \emph{expectation} of the residual
quadratic form rather than its plug-in value.  Under the Gaussian
approximation to the state posterior, that expectation is the plug-in
value plus a covariance correction,
\begin{equation}
\E\bigl[\,r_t^\top R^{-1}r_t \,\big|\, Y\bigr]
= \widehat r_t^{\,\top} R^{-1}\widehat r_t
  + \tr\!\bigl(R^{-1}H\,\Sigma_{tt}\,H'\bigr),
\label{eq:exact-estep-meas}
\end{equation}
with $\Sigma_{tt}$ delivered by the selected inversion of
Section~\ref{sec:selinv}.  For Student-$t$ the package implements this
covariance-corrected E-step as an option
(\texttt{opts.exact\_estep} in Octave, \texttt{exact\_estep=True} in
Python), at the cost of one selected inversion per iteration; the
corrected weights are smaller in periods of high state uncertainty,
which is the honest behaviour --- an apparently clean residual at an
imprecisely estimated state is weaker evidence against an outlier than
the plug-in update assumes.

\subsubsection{State-side robustness}

The same construction applies to the state innovations $w_t$.  Heavy
tails on $w_t$ accommodate occasional large jumps in the latent path
while still smoothing tightly through quiet periods.  The IRLS weights
$\xi_t$ on $w_t$ are updated by exactly the same formulas as $\tau_t$,
using the squared Mahalanobis distance (the quadratic form
$\widehat w_t^\top Q^{-1}\widehat w_t$, not its square root) of the
smoothed innovation
\begin{equation}
\widehat w_t = \widehat x_{t\mid T} - F \widehat x_{t-1\mid T} - Bz_t,
\end{equation}
and the time-varying state covariance becomes $Q_t = Q/\xi_t$.  The
exogenous forcing $Bz_t$ must be subtracted here for the same reason
as on the measurement side: otherwise known exogenous movements in the
state would be misread as structural breaks.  The covariance-corrected
E-step analogue of \eqref{eq:exact-estep-meas} uses
\begin{equation}
\E\bigl[\,w_t^\top Q^{-1}w_t \,\big|\, Y\bigr]
= \widehat w_t^{\,\top} Q^{-1}\widehat w_t
 + \tr\!\Bigl(Q^{-1}\bigl[\Sigma_{tt}-F\Sigma_{t,t-1}'-\Sigma_{t,t-1}F'
 +F\Sigma_{t-1,t-1}F'\bigr]\Bigr),
\end{equation}
which requires the lag-one block $\Sigma_{t,t-1}$ in addition to the
diagonal blocks --- both already produced by the same Takahashi sweep.

This is the practical machinery for structural-break detection.  Set
the state-innovation prior to Student-$t$ with $\nu = 4$, run the
smoother, and read off the periods at which $\xi_t \ll 1$:  those are
the breaks.  The demo \texttt{demo\_robust\_breaks.m} does this on a
near-random-walk series with three injected level shifts; the
algorithm identifies all three at $\xi_t < 0.02$ while
$\xi_t \approx 1$ elsewhere, and the smoothed path achieves under a
third of the RMSE of the Gaussian smoother.

\subsubsection{Implementation}

The package implements all of the above in
\texttt{simks\_smooth\_robust.m}.  The interface is
\begin{lstlisting}[language=Octave]
[Xhat, info] = simks_smooth_robust(Y, F, H, Q, R, a0, P0, family, opts);
\end{lstlisting}
with \texttt{family} one of \texttt{'gaussian'}, \texttt{'student-t'},
\texttt{'laplace'}, \texttt{'huber'}, and \texttt{opts.robust\_state}
optionally activating the state-side variant.  The returned
\texttt{info.tau\_v} (and \texttt{info.tau\_w}) is the per-period
weight vector at convergence; values much smaller than $1$ are
diagnostic flags for outlier observations or structural breaks.

The demo \texttt{demo\_robust.m} contaminates 10\% of observations
with extra noise and compares the Gaussian, Student-$t$, Huber, and
Laplace smoothers.  Student-$t$ and Huber cut RMSE by roughly
40~percent versus Gaussian and correctly downweight most outlier
observations.

\subsubsection{Beyond scale mixtures}

The IRLS-on-the-same-$J$ scheme covers more than Gaussian scale
mixtures.  Asymmetric Laplace gives quantile-style smoothing; smooth
approximations of the $\ell_1$ penalty give pseudo-Huber; generalized
linear measurement models (Poisson, binomial, multinomial) give
non-Gaussian observation links, with the IRLS weights determined by
the local curvature of the link function.  All of these inherit the
sparsity structure of $J$ automatically.

Thus the simultaneous view is not merely an alternative derivation of
the Kalman smoother.  It is the right abstraction for a broader class
of sparse latent-path estimation problems; see \citet{RueHeld2005} for
the spatial-statistics counterpart.

\subsection{Regime-switching dynamics}
\label{sec:slds}

Macroeconomic models often need state-space dynamics that depend on a
discrete unobserved regime --- expansion versus recession, dovish
versus hawkish monetary policy, high versus low financial-stress
periods.  The natural model is a switching linear dynamical system
(SLDS):
\begin{align}
x_t &= F_{s_t}\,x_{t-1}+B_{s_t}\,z_t+w_t,
        \quad w_t\sim N(0,Q_{s_t}),\\
y_t &= H x_t+D z_t+v_t,\quad v_t\sim N(0,R),
\end{align}
with $s_t\in\{1,\ldots,K\}$ a Markov chain on $K$ regimes with
transition matrix $A$.

The exact posterior $p(X,S\mid Y)$ is a Gaussian mixture with $K^T$
components and intractable for any non-trivial $T$.  The standard
scalable approach, due to \citet{GhahramaniHinton2000}, is variational
EM with a single-Gaussian-times-Markov-chain factorization:
\begin{equation}
q(X,S)=q(X)\prod_{t=1}^T q(s_t).
\end{equation}
This decouples the two latent variables and turns inference into
alternating updates of two tractable problems.

\subsubsection{Variational EM iteration}

Let $\pi_{t,k}=q(s_t=k)$ be the regime responsibilities.  Coordinate
ascent on the evidence lower bound alternates two exact updates, and
it is worth being careful about what ``exact'' means here, because a
tempting shortcut --- smoothing under the responsibility-\emph{averaged}
matrices $\widetilde F_t=\sum_k\pi_{t,k}F_k$,
$\widetilde Q_t=\sum_k\pi_{t,k}Q_k$ --- is a moment-matching
approximation, \emph{not} the variational update, and the two differ
whenever the regimes are uncertain (note in particular that
$\sum_k\pi_{t,k}Q_k^{-1}\neq(\sum_k\pi_{t,k}Q_k)^{-1}$).  The package
implements the exact updates.
\begin{enumerate}[leftmargin=2em]
\item \textbf{$q(X)$-step (smooth under expected quadratic forms).}
  The optimal coordinate update sets
  $\log q(X)=\E_{q(S)}\log p(X,S,Y)+\text{const}$, which is Gaussian,
  with the expected transition term at $t$ averaging the per-regime
  quadratic forms.  Collecting terms in $X$, the period-$t$
  contributions to the precision matrix are
  \begin{equation}
  \overline{Q^{-1}}_t=\sum_{k=1}^K \pi_{t,k}Q_k^{-1},
  \quad
  \overline{Q^{-1}F}_t=\sum_{k=1}^K \pi_{t,k}Q_k^{-1}F_k,
  \quad
  \overline{F'Q^{-1}F}_t=\sum_{k=1}^K \pi_{t,k}F_k'Q_k^{-1}F_k,
  \end{equation}
  assembled as
  $J_{tt}=\overline{Q^{-1}}_t+\overline{F'Q^{-1}F}_{t+1}+H'R^{-1}H$
  and $J_{t,t-1}=-\overline{Q^{-1}F}_t$ (with the obvious endpoint
  modifications), and with the information vector carrying the
  matching mixed exogenous terms
  $\overline{Q^{-1}B}_t\,z_t$ and $-\overline{F'Q^{-1}B}_{t+1}\,z_{t+1}$.
  Crucially, this precision matrix is \emph{still block tridiagonal at
  every iteration} --- the expected quadratic forms couple only
  adjacent periods, exactly as in the fixed-regime model --- so each
  $q(X)$-step is one sparse Cholesky of the usual structure; only the
  block values change.
\item \textbf{$q(S)$-step (regime update under expected
  log-likelihoods).}  The optimal $q(s_t)$ update requires the
  \emph{expected} Gaussian transition log-likelihood under $q(X)$,
  not the density evaluated at the smoothed point path.  Writing
  $\widehat d_{t,k}=\widehat x_t-F_k\widehat x_{t-1}-B_kz_t$,
  \begin{equation}
  \begin{split}
  \ell_{t,k}
  &=\E_{q(X)}\log N\bigl(x_t;F_kx_{t-1}+B_kz_t,Q_k\bigr)\\
  &=-\tfrac12\Bigl[\widehat d_{t,k}^{\,\top}Q_k^{-1}\widehat d_{t,k}
   +\tr\bigl(Q_k^{-1}\mathcal C_{t,k}\bigr)\Bigr]
   -\tfrac12\log\det Q_k-\tfrac n2\log 2\pi,
  \end{split}
  \end{equation}
  with the covariance correction
  \begin{equation}
  \mathcal C_{t,k}
  =\Sigma_{tt}-F_k\Sigma_{t,t-1}'-\Sigma_{t,t-1}F_k'
   +F_k\Sigma_{t-1,t-1}F_k'
  \end{equation}
  built from the diagonal and lag-one blocks that the selected
  inversion of Section~\ref{sec:selinv} already produces.  The
  standard discrete forward--backward HMM smoother is then run on the
  pair $(A,\,e^{\ell})$, at $O(TK^2)$ cost, dominated by the state
  smoothing.
\end{enumerate}
The two steps are iterated to convergence in $\pi$.  Typically 10--40
iterations suffice for moderate-$T$ macroeconomic data.  The
covariance corrections matter in both directions: omitting the trace
terms (i.e.\ scoring regimes at the point path) systematically
overstates regime-classification confidence, because a residual that
looks small at an imprecisely estimated state is weaker evidence than
the same residual at a well-pinned-down state.  A corollary worth
knowing when designing a regime-switching exercise: regimes that
differ mainly in innovation \emph{variance} are sharply classifiable
only when the posterior state uncertainty is small relative to the
quiet regime's innovation variance --- otherwise even an oracle
observing the true states classifies imperfectly.

The substantive payoff is that the precision-form formulation extends
to regime-switching dynamics without any new mathematical machinery:
every $q(X)$-step is one sparse Cholesky on a block-tridiagonal $J$ of
the same pattern the no-switching smoother uses, and the covariance
blocks the $q(S)$-step needs come from the same Takahashi sweep that
EM already requires.  The recursive form, by contrast, loses its
closed-form structure once regimes are introduced ---
exact filtering would require maintaining $K^t$ Gaussians at time $t$,
with the Kim and Hamilton filters performing per-step collapses to a
single Gaussian per regime.  The simultaneous form makes a
single-Gaussian approximation too, but as the optimal $q(X)$ within
the variational family, with a monotone evidence-lower-bound
interpretation rather than an ad hoc collapse.

The demo \texttt{demo\_slds.m} implements a two-regime AR(1) with
expansion/recession dynamics (persistence $\phi_1=0.85$,
$\phi_2=0.50$; volatility $\sigma^2_1=0.02$, $\sigma^2_2=0.20$;
measurement noise $R=0.01$, small relative to $\sigma_1^2$ per the
corollary above).  The demo generates its data with a small portable
random-number generator (a Park--Miller linear congruential generator
plus Box--Muller, implemented identically in the Octave/MATLAB and
Python versions), so all three environments simulate the same data
and print the same numbers: regime classification accuracy
$\approx 0.90$, mean $\pi_{t,\text{rec}}\approx 0.80$ during true
recessions and $\approx 0.06$ during true expansions, and SLDS state
RMSE of $0.095$ against $0.099$ for the single-regime Gaussian
smoother run at averaged dynamics.

\subsection{Constrained smoothing}
\label{sec:constrained}

Many macroeconomic state-space objects have known structural bounds
that the textbook Gaussian smoother ignores.  A NAIRU estimate that
goes negative is meaningless; a debt-to-GDP latent component that
exceeds historical maxima is implausible; an equilibrium real interest
rate is sometimes assumed monotone over a regime.  These are
inequality constraints on the smoothed path,
\begin{equation}
A_{\text{ineq}}\,X\le b_{\text{ineq}}.
\label{eq:ineq}
\end{equation}
Such constraints have no clean recursive-filter counterpart; the
Kalman recursion produces a Gaussian conditional, which is in general
not respected by any non-trivial inequality.

\subsubsection{Constrained smoothing as a sparse QP}

In the simultaneous formulation, constrained smoothing is the
quadratic program
\begin{equation}
\min_X\ \tfrac12 X'JX-h'X
\quad\text{s.t.}\quad A_{\text{ineq}}X\le b_{\text{ineq}},
\label{eq:qp}
\end{equation}
with $J$ and $h$ exactly as before.  When the constraints are
inactive at the unconstrained optimum, $\widehat X=J^{-1}h$ already
solves~\eqref{eq:qp}; when some constraints bind, the solution is
the constrained posterior mode.

\subsubsection{Interior-point algorithm on the same sparse $J$}

The standard primal-dual interior-point method uses a log-barrier:
\begin{equation}
X(\mu)=\argmin_X\ \tfrac12 X'JX-h'X-\mu\sum_{i=1}^{c}\log(b_i-a_i'X).
\end{equation}
The Newton step at $X$ requires solving
\begin{equation}
\Bigl(J+\mu\,A_{\text{ineq}}^\top\diag\bigl(1/(b-A_{\text{ineq}}X)^2\bigr)A_{\text{ineq}}\Bigr)\,\Delta X
=-\nabla,
\label{eq:newton}
\end{equation}
where the gradient is
$\nabla=JX-h+\mu A_{\text{ineq}}^\top(1/(b-A_{\text{ineq}}X))$.
The sign of the barrier term is worth a sanity check, since the
displayed Hessian only pins it down implicitly: differentiating
$-\mu\log(b_i-a_i'X)$ gives $+\mu\,a_i/(b_i-a_i'X)$, a force that
blows up as the iterate approaches the boundary and pushes the
minimizer back into the interior, exactly as a barrier should; with
the opposite sign the ``barrier'' would attract iterates to the
boundary and the method would degenerate into clipping.

The structural point is that the Hessian is $J$ plus a low-rank
positive-semidefinite update.  For typical macroeconomic
constraints --- bounds on a single state coordinate at each $t$, or
on adjacent differences --- $A_{\text{ineq}}$ is itself extremely
sparse, often with only $O(n)$ nonzeros per row.  The augmented
Hessian $J+\mu A_{\text{ineq}}^\top D\,A_{\text{ineq}}$ then remains
block tridiagonal (component-wise bounds) or grows the bandwidth by
one block (monotonicity constraints).  Either way, sparse Cholesky
still applies, and the per-iteration cost stays $O(Tn^3)$.

The barrier parameter $\mu$ is decreased geometrically until the
duality gap is below tolerance.  Total cost is the number of Newton
iterations times the sparse-Cholesky cost; typical convergence is
20--60 inner iterations across 30--40 barrier reductions for a
moderate-sized state-space problem with component-wise bounds.

\subsubsection{Posterior interpretation}

The QP solution is the constrained posterior \emph{mode}, not the
posterior mean.  Under truncated Gaussian posteriors the mean differs
from the mode in general; computing the mean requires sampling
methods such as \citet{Botev2017}'s exact truncated-Gaussian
algorithm.  For most applied macro uses --- displaying a smoothed
trend that respects non-negativity, or a credible band whose
endpoints respect a known structural bound --- the mode is the
quantity of interest.

The demo \texttt{demo\_constrained.m} imposes $x_t\ge 0$ on a noisy
near-zero AR(1) smoother (a stand-in for a NAIRU-like object) and
verifies that the constraint is exactly satisfied while the
unconstrained smoother places $31/120$ periods below zero.  RMSE
against the true non-negative trend improves by $\sim 4\%$, illustrating
that imposing known structural information is information-efficient,
not just cosmetic.

\subsubsection{Inequality constraints beyond bounds}

The same QP machinery handles a broad set of constraints:
\begin{itemize}[leftmargin=2em]
\item \textbf{Box constraints.}  $\ell_t\le x_{t,j}\le u_t$ on
  selected coordinates and periods: two rows per constraint.
\item \textbf{Monotonicity.}  $x_{t+1,j}\ge x_{t,j}$ for a trend
  component: rows that pick out $-x_{t+1,j}+x_{t,j}$.
\item \textbf{Smoothness / curvature.}  Bounds on
  $\Delta^2 x_t=x_t-2x_{t-1}+x_{t-2}$: three-period stencils.
\item \textbf{Cone constraints.}  Bounds on $\|x_t\|_\infty$ etc.,
  representable as collections of linear inequalities.
\end{itemize}
The simultaneous form turns each into a few extra sparse rows of
$A_{\text{ineq}}$, with no impact on the underlying $J$.  This is
the kind of extension that is essentially impossible to graft onto a
recursive Kalman filter and easy to add here.

\section{What you gain}

Beyond working in the right idiom, the simultaneous formulation has
four concrete advantages over the recursive form.
\begin{enumerate}[leftmargin=2em]
\item \textbf{Numerical robustness, properly scoped.}  Sparse Cholesky
  on $J$ maintains symmetry and positive definiteness by construction:
  $J$ is assembled, not propagated, so the failure mode of a hand-coded
  recursive filter --- losing definiteness through subtractions of
  large covariance matrices, the reason the square-root, UDU, and
  Bierman--Thornton variants exist --- cannot occur.  The honest
  counterpart, spelled out in Section~\ref{sec:where-used}, is that the
  precision form is a normal-equations representation whose conditioning
  degrades quadratically where square-root forms degrade linearly, and
  which needs SPD $Q$ and $R$ (or explicit regularization).  For
  well-posed batch smoothing problems --- the use cases of this note ---
  the Cholesky route is robust and simple; for genuinely
  ill-conditioned online problems the square-root literature retains
  its documented advantages.
\item \textbf{Modularity.}  Missing data, diffuse priors, irregular
  sampling, partial observations, time-varying parameters, and
  unobserved-components decompositions are all small edits to the
  assembly of $J$ and $h$.  No need to redesign a recursion or
  re-derive boundary conditions.
\item \textbf{Extensibility.}  Robust losses (Laplace, Student-$t$),
  sparsity priors, inequality constraints, and hierarchical extensions
  all plug into the same sparse linear-algebra framework, typically via
  iteratively reweighted solves on the same $J$.
\item \textbf{Conceptual unification.}  The Hodrick--Prescott filter,
  the Beveridge--Nelson decomposition, unobserved-components models,
  factor models, and the textbook state-space problem are all
  instances of a single object: a quadratic penalty on a latent path
  solved by $J\widehat X=h$.  This unification is hard to see through
  the recursive lens.
\end{enumerate}

What you do not gain is raw speed.  Sparse Cholesky on a
block-tridiagonal SPD matrix \emph{is} the block-LDL$^\top$ recursion;
CHOLMOD finds the natural ordering for free, because the time ordering
is already optimal.  Both methods are $O(Tn^3)$ and the constants are
comparable.  The point of the simultaneous form is not that it computes
the answer faster.  It is that it computes the answer in a form the
student can recognize as the answer.

\section{Conclusion}

The recursive Kalman filter and the RTS smoother are not separate
theories.  They are the block-LDL$^\top$ forward sweep and the
selected-inversion backward sweep of one sparse block-tridiagonal
linear system,
\begin{equation*}
J\widehat X=h.
\end{equation*}
The reason this is not how they are usually taught is historical: the
recursive form was forced on Kalman and his contemporaries by 1960s
memory constraints and shaped by an engineering rather than an
econometric tradition.  Neither constraint is ours.  For batch
state-space estimation in modern macroeconomics, the simultaneous form
is in the native idiom; numerically robust for well-posed problems with
SPD precision, with none of the subtraction-induced failure modes of
hand-coded recursions; modular under missing data and diffuse priors;
easily extended to non-Gaussian errors and unobserved-components
decompositions; and at least as efficient.
Posterior covariances, marginal likelihood, and EM with smoothed
sufficient statistics all follow from the same sparse solve.

\clearpage
\appendix

\section*{License}
\addcontentsline{toc}{section}{License}

Copyright \copyright\ 2026 Mico Mrkaic.  The package is released under
the GNU General Public License version~3 or later
(\url{https://www.gnu.org/licenses/gpl-3.0.html}).  In plain
terms: you may use, modify, and redistribute the code freely, including
for commercial purposes, provided that derivative works are released
under the same license with source code available.  Every source file
in the package carries a short GPLv3 header pointing to the full
license text in the \texttt{LICENSE} file at the package root.

\section{Octave reference implementation}
\label{sec:octave-impl}

The reference implementation lives in \texttt{octave/simksmoother/}
and is written in pure Octave (also runs on MATLAB).  It assembles
$J$ and $h$ directly as sparse matrices via triplet arrays $(i,j,v)$
with one terminal call to \texttt{sparse(i,j,v,N,N)}.  No matter which
routine is invoked, the fundamental operation is one sparse solve:
\begin{lstlisting}[language=Octave]
Xvec = J \ h;
\end{lstlisting}
Octave's backslash dispatches to CHOLMOD when $J$ is sparse and SPD,
so the cost is $O(Tn^3)$ in both flops and memory, just as in the
recursive Kalman filter.

\subsection{Calling conventions}

All routines share the same calling style.  Required arguments
(\texttt{Y}, \texttt{F}, \texttt{H}, \texttt{Q}, \texttt{R},
\texttt{a0}, \texttt{P0}) come first; optional behaviour is controlled
through an \texttt{opts} struct passed as the final argument.  Common
fields:
\begin{itemize}[leftmargin=2em,itemsep=2pt]
\item \texttt{opts.Z}\ : $k\times T$ exogenous inputs.  Omit for none.
\item \texttt{opts.B}\ : $n\times k$ state-side loading (default zeros).
\item \texttt{opts.D}\ : $m\times k$ measurement-side loading (default zeros).
\item \texttt{opts.return\_cov}\ : if \texttt{true}, also return posterior covariance blocks via selected inversion.
\end{itemize}
Diffuse priors are signalled by \texttt{P0 = []}.  Missing
observations are \texttt{NaN} entries in \texttt{Y}; they are dropped
from $\mathcal H^\top R^{-1}\mathcal H$ at the affected periods
automatically.  The path \texttt{Xhat} returned by every routine is
$n\times(T+1)$, with columns $\widehat x_{0\mid T},\ldots,
\widehat x_{T\mid T}$.

\subsection{Core smoothing and inference}

\paragraph{\texttt{simks\_smooth\_const}.}  Basic constant-parameter
smoother; the workhorse.
\begin{lstlisting}[language=Octave]
[Xhat, J, h, info] = simks_smooth_const(Y, F, H, Q, R, a0, P0, opts)
\end{lstlisting}
\emph{Inputs:} the seven required state-space arguments plus the
optional \texttt{opts} struct.  \emph{Outputs:} the smoothed path
\texttt{Xhat}; the sparse precision matrix \texttt{J}; the
information vector \texttt{h}; and an \texttt{info} struct that
always contains $n$, $m$, $T$ and the nonzero count of $J$, and
additionally \texttt{Ptt}, \texttt{Ptt1} (cell arrays of diagonal and
first-subdiagonal blocks of $J^{-1}$) when
\texttt{opts.return\_cov=true}.

\paragraph{\texttt{simks\_smooth\_tv}.}  Time-varying smoother.
\begin{lstlisting}[language=Octave]
[Xhat, J, h, info] = simks_smooth_tv(Y, Fseq, Hseq, Qseq, Rseq, a0, P0, opts)
\end{lstlisting}
Each of \texttt{Fseq}, \texttt{Hseq}, \texttt{Qseq}, \texttt{Rseq}
may be a single constant matrix (replicated across all $t$) or a
$1\times T$ cell array of per-period matrices; mixed usage is
supported.  Outputs as in \texttt{simks\_smooth\_const}.

\paragraph{\texttt{simks\_selected\_inv}.}  Takahashi selected inversion.
\begin{lstlisting}[language=Octave]
[Ptt, Ptt1, logdetJ] = simks_selected_inv(J, n)
\end{lstlisting}
\emph{Inputs:} the precision matrix \texttt{J} and the block size
$n$.  \emph{Outputs:} the diagonal blocks $\Sigma_{tt}$ and the
first-subdiagonal blocks $\Sigma_{t,t-1}$ of $J^{-1}$ as cell arrays,
plus $\log\det J$ as a byproduct.  These are precisely the
sufficient statistics the EM E-step needs.

\paragraph{\texttt{simks\_marginal\_loglik}.}  Marginal log-likelihood.
\begin{lstlisting}[language=Octave]
[ll, parts] = simks_marginal_loglik(Y, F, H, Q, R, a0, P0, opts)
\end{lstlisting}
Returns the scalar marginal log-likelihood $\log p(Y;\theta)$ via the
formula in Section~\ref{sec:loglik}.  The \texttt{parts} struct
breaks out the prior, data, $\log\det J$, and normalizing
contributions for diagnostics.

\paragraph{\texttt{simks\_simulate\_const}.}  Forward simulation.
\begin{lstlisting}[language=Octave]
[Y, X] = simks_simulate_const(T, F, H, Q, R, a0, P0, opts)
\end{lstlisting}
\emph{Inputs:} the horizon $T$ and the state-space parameters.
\emph{Outputs:} \texttt{Y} ($m\times T$) and \texttt{X}
($n\times(T+1)$, columns $x_0,\ldots,x_T$).  Used by the demos and
the test suite to generate ground-truth data.

\paragraph{\texttt{simks\_rts\_reference}.}  Textbook Kalman + RTS.
\begin{lstlisting}[language=Octave]
[Xhat, P] = simks_rts_reference(Y, F, H, Q, R, a0, P0)
\end{lstlisting}
A clean recursive implementation, provided so that the reader can
compare the recursive and simultaneous formulations line by line and
confirm they compute the same object to machine precision.  No
exogenous-input or missing-data handling --- it is the textbook
algorithm in its simplest form.

\subsection{Parameter estimation}

\paragraph{\texttt{simks\_em\_const}.}  Gaussian EM.
\begin{lstlisting}[language=Octave]
[params, Xhat, hist] = simks_em_const(Y, n, opts)
\end{lstlisting}
\emph{Inputs:} the data \texttt{Y}, the latent state dimension $n$,
and an \texttt{opts} struct that may carry initial parameter values,
convergence controls, and exogenous inputs \texttt{Z}.
\emph{Outputs:} the estimated parameter struct \texttt{params} with
fields \texttt{F}, \texttt{H}, \texttt{Q}, \texttt{R}, \texttt{a0},
\texttt{P0} (plus \texttt{B}, \texttt{D} if exogenous was supplied);
the final smoothed path; and the history of the marginal
log-likelihood across iterations.  The history is monotone
non-decreasing by construction and is a useful sanity check.  All
returned parameters, including \texttt{a0} and \texttt{P0}, are
estimated: each iteration updates the initial-state pair by the
single-observation M-step $\widehat a_0=\widehat x_{0\mid T}$ and
$\widehat P_0=\Sigma_{00}$ documented in Section~\ref{sec:em}.

\subsection{Beyond-Gaussian and constrained smoothing}

\paragraph{\texttt{simks\_smooth\_robust}.}  Heavy-tailed IRLS smoother.
\begin{lstlisting}[language=Octave]
[Xhat, info] = simks_smooth_robust(Y, F, H, Q, R, a0, P0, family, opts)
\end{lstlisting}
\emph{Inputs:} the standard state-space arguments plus the
\texttt{family} string and \texttt{opts}.  The family is one of
\texttt{'gaussian'}, \texttt{'student-t'}, \texttt{'laplace'}, or
\texttt{'huber'}.  The
\texttt{opts} struct may carry \texttt{nu} (Student-$t$ degrees of
freedom), \texttt{c} (Huber tuning constant), \texttt{robust\_state}
(boolean, applies heavy tails to $w_t$ as well),
\texttt{exact\_estep} (boolean, Student-$t$ only: replaces plug-in
IRLS with the covariance-corrected E-step of
Section~\ref{sec:robust}), and the usual exogenous fields.
\emph{Outputs:} the smoothed path and an \texttt{info} struct that
always contains \texttt{tau\_v}
(per-period measurement weights at convergence), and additionally
\texttt{tau\_w} when \texttt{robust\_state} is \texttt{true}.
Weights much below $1$ flag outliers or structural breaks.

\paragraph{\texttt{simks\_smooth\_slds}.}  Switching dynamics.
\begin{lstlisting}[language=Octave]
[Xhat, info] = simks_smooth_slds(Y, F_set, H, Q_set, R, a0, P0, opts)
\end{lstlisting}
Variational EM (exact coordinate updates, Section~\ref{sec:slds})
for the switching linear dynamical system.  \emph{Inputs:}
\texttt{F\_set} and \texttt{Q\_set} are $1\times K$ cell arrays of
per-regime matrices; \texttt{opts.A} is the $K\times K$
regime-transition matrix; with exogenous inputs, the per-regime state
loadings are passed as \texttt{opts.B\_set}, a $1\times K$ cell
array of $n\times k$ matrices (defaulting to zeros), alongside the
common \texttt{opts.Z} and the regime-invariant \texttt{opts.D} ---
note that the \texttt{opts.B} field of \texttt{simks\_smooth\_const}
is \emph{not} used here, precisely because the loading is
regime-dependent.  The remaining arguments are as in
\texttt{simks\_smooth\_const}.  \emph{Outputs:} the smoothed state
path; \texttt{info.pi}, a $K\times T$ matrix of regime
responsibilities $\pi_{t,k}$; and \texttt{info.Ptt},
\texttt{info.Ptt1}, the posterior covariance blocks of $q(X)$ at
convergence (used internally by the expected-log-likelihood regime
update and returned for convenience).

\paragraph{\texttt{simks\_smooth\_constrained}.}  Inequality-constrained smoother.
\begin{lstlisting}[language=Octave]
[Xhat, info] = simks_smooth_constrained(Y, F, H, Q, R, a0, P0, A_ineq, b_ineq, opts)
\end{lstlisting}
\emph{Inputs:} the standard state-space arguments plus the constraint
data $A_{\text{ineq}}$ (sparse, $c\times(T+1)n$) and $b_{\text{ineq}}$
(length $c$).  \emph{Outputs:} the constrained smoothed path
\texttt{Xhat} and an \texttt{info} struct containing the active-set
mask, the approximate KKT multipliers, and interior-point convergence
diagnostics.

\subsection{Demos}

The \texttt{octave/examples/} directory contains eleven demos
exercising every routine and feature of the package.  Every demo
prints numerical diagnostics to the console and produces a figure
showing the relevant smoothed trajectories, posterior bands, IRLS
weights, or regime probabilities.  Figures are simultaneously saved to
\texttt{octave/examples/figures/} and displayed on screen.  The three
worked applications in Appendix~\ref{sec:apps} produce figures in
the same way.  The driver script \texttt{demo\_main.m} presents an
interactive menu and runs demos individually or in sequence:
\begin{lstlisting}[language=Octave]
cd octave/examples
demo_main
\end{lstlisting}
The demos are:
\begin{itemize}[leftmargin=2em,itemsep=2pt]
\item \texttt{demo\_const\_known}: smoothing with known parameters; posterior bands and empirical coverage check on a 2D state.
\item \texttt{demo\_tv\_known}: time-varying parameters via the cell-array interface; a sinusoidally varying $F_t$ illustrates the API.
\item \texttt{demo\_compare\_recursive}: simultaneous smoother vs textbook Kalman + RTS, agreement to $\sim 10^{-15}$.
\item \texttt{demo\_missing}: NaN observations dropped automatically; posterior variance grows inside the gap.
\item \texttt{demo\_hp\_filter}: the HP filter as a special case, matching the closed-form $(I+\lambda D_2^\top D_2)^{-1}y$ to within boundary effects.
\item \texttt{demo\_const\_em}: EM with smoothed sufficient statistics; verifies that the marginal log-likelihood is monotone non-decreasing.
\item \texttt{demo\_robust}: 10\% of observations contaminated with extra noise; Student-$t$ and Huber cut RMSE by roughly 40~percent.
\item \texttt{demo\_robust\_breaks}: structural-break detection via state-side robust prior; three injected level shifts identified at $\xi_t<0.02$.
\item \texttt{demo\_hp\_credit}: flagship demo on an HP filter with a credit-cycle exogenous driver, breaks, and outliers; compares Gaussian smoothing, Student-$t$ robust smoothing at true parameters, and Gaussian EM for unknown $(F,H,Q,R,B,D)$ as separate estimators (see Section~\ref{sec:em}).
\item \texttt{demo\_slds}: two-regime AR(1); exact variational EM recovers regime probabilities at $\sim 90\%$ classification accuracy.  This demo uses a portable random-number generator so Octave, MATLAB, and Python simulate identical data and report identical results.
\item \texttt{demo\_constrained}: NAIRU-style non-negativity constraint via interior-point QP; constraint exactly satisfied while the unconstrained smoother goes negative in 31 of 120 periods.
\end{itemize}
The three worked applications of Appendix~\ref{sec:apps} are also
implemented as runnable scripts, with filenames \texttt{app$k$\_*.m},
in the \texttt{applications/} subdirectory of
\texttt{octave/examples/}.

\section{Python implementation}
\label{sec:python-impl}

The \texttt{python/} directory contains a direct port of the Octave
implementation to Python, packaged as \texttt{simksmoother}.  The
mathematics is identical to the Octave version; only surface
conventions differ.

\subsection{Installation}

The recommended setup is the standard editable install:
\begin{lstlisting}
cd python
pip install -e .
\end{lstlisting}
Requires Python 3.9 or later, NumPy 1.20+, and SciPy 1.7+, plus
\texttt{matplotlib} for the demo figures.  Optional:
\texttt{scikit-sparse} (CHOLMOD bindings) for faster sparse
factorization on large problems.

The demos and worked applications also run without
\texttt{pip install -e .}\ by adding the \texttt{python/} directory
to \texttt{sys.path} themselves when they detect that
\texttt{simksmoother} is not yet importable.  This lets users explore
the package straight from a freshly unzipped copy of the repository
(e.g.\ via IPython's \texttt{run demo\_const\_known} from inside
\texttt{python/examples/}).

\subsection{Calling conventions}

Functions are imported from the top-level package:
\begin{lstlisting}[language=Python]
import simksmoother as sks
\end{lstlisting}
Required arguments come first; options are passed as keyword
arguments rather than an \texttt{opts} struct.  \texttt{Z},
\texttt{B}, \texttt{D} keyword arguments carry exogenous inputs.
Diffuse priors are signalled by \texttt{P0=None} (instead of Octave's
\texttt{P0=[]}).  Missing observations are \texttt{np.nan} entries in
\texttt{Y}.  Returned arrays use the NumPy convention: rows are
state coordinates, columns are time.  Note the indexing offset
relative to Octave: column \texttt{t} of \texttt{Xhat} corresponds
to $x_t$ ($t = 0,1,\ldots,T$), not column \texttt{t+1}.

\subsection{Core smoothing and inference}

\paragraph{\texttt{smooth\_const}.}  Basic smoother.
\begin{lstlisting}[language=Python]
Xhat, J, h, info = sks.smooth_const(
    Y, F, H, Q, R, a0, P0,
    Z=None, B=None, D=None, return_cov=False)
\end{lstlisting}
\emph{Inputs:} the seven required state-space arguments; optional
exogenous inputs \texttt{Z}, \texttt{B}, \texttt{D}; flag
\texttt{return\_cov}.  \emph{Outputs:} the smoothed path
\texttt{Xhat} ($n\times(T+1)$ NumPy array), the SciPy sparse
precision matrix \texttt{J} (CSC format), the information vector
\texttt{h}, and \texttt{info} (a Python dict).  With
\texttt{return\_cov=True}, \texttt{info} contains
\texttt{Ptt} and \texttt{Ptt1}, lists of NumPy arrays for the
diagonal and first-subdiagonal blocks of $J^{-1}$.

\paragraph{\texttt{smooth\_tv}.}  Time-varying smoother.
\begin{lstlisting}[language=Python]
Xhat, J, h, info = sks.smooth_tv(
    Y, Fseq, Hseq, Qseq, Rseq, a0, P0,
    Z=None, Bseq=None, Dseq=None, return_cov=False)
\end{lstlisting}
Each of \texttt{Fseq}, \texttt{Hseq}, \texttt{Qseq}, \texttt{Rseq}
may be either a single NumPy array (constant) or a length-$T$ list
of arrays (per-period).

\paragraph{\texttt{selected\_inv}.}  Selected inversion.
\begin{lstlisting}[language=Python]
Ptt, Ptt1 = sks.selected_inv(J, n)
\end{lstlisting}
\emph{Inputs:} the sparse precision matrix \texttt{J} and the block
size $n$.  \emph{Outputs:} Python lists of $(n,n)$ NumPy arrays for
the diagonal blocks (length $T+1$) and the first-subdiagonal blocks
(length $T$) of $J^{-1}$.

\paragraph{\texttt{marginal\_loglik}.}  Marginal log-likelihood.
\begin{lstlisting}[language=Python]
ll, parts = sks.marginal_loglik(
    Y, F, H, Q, R, a0, P0, Z=None, B=None, D=None)
\end{lstlisting}
\emph{Outputs:} the scalar log-likelihood \texttt{ll} and a dict
\texttt{parts} with the diagnostic breakdown.

\paragraph{\texttt{simulate\_const}.}  Forward simulation.
\begin{lstlisting}[language=Python]
Y, X = sks.simulate_const(
    T, F, H, Q, R, a0, P0,
    Z=None, B=None, D=None, rng=None)
\end{lstlisting}
Uses the modern \texttt{numpy.random.Generator} API; pass an
\texttt{rng} for reproducibility or accept the default.

\paragraph{\texttt{rts\_smoother}.}  Reference Kalman + RTS.
\begin{lstlisting}[language=Python]
Xhat, P = sks.rts_smoother(Y, F, H, Q, R, a0, P0)
\end{lstlisting}

\subsection{Parameter estimation}

\paragraph{\texttt{em\_const}.}  Gaussian EM.
\begin{lstlisting}[language=Python]
params, Xhat, hist = sks.em_const(
    Y, n, max_iter=200, tol=1e-6,
    F=None, H=None, Q=None, R=None, a0=None, P0=None,
    Z=None, B=None, D=None, verbose=False)
\end{lstlisting}
\emph{Inputs:} the data and the state dimension $n$; optional
initializations and convergence controls; optional exogenous inputs.
\emph{Outputs:} \texttt{params} (dict with keys \texttt{F},
\texttt{H}, \texttt{Q}, \texttt{R}, \texttt{a0}, \texttt{P0}, and
\texttt{B}, \texttt{D} if exogenous was supplied), the final smoothed
path, and a list of marginal log-likelihood values across iterations
(monotone non-decreasing).

\subsection{Beyond-Gaussian and constrained smoothing}

\paragraph{\texttt{smooth\_robust}.}  Heavy-tailed IRLS smoother.
\begin{lstlisting}[language=Python]
Xhat, info = sks.smooth_robust(
    Y, F, H, Q, R, a0, P0, family,
    nu=4.0, c=1.345, robust_state=False,
    exact_estep=False,
    max_iter=50, tol=1e-6,
    Z=None, B=None, D=None)
\end{lstlisting}
\texttt{family} is one of \texttt{"gaussian"}, \texttt{"student-t"},
\texttt{"laplace"}, \texttt{"huber"}.  By default the weights are
plug-in IRLS (the fixed point is the posterior mode);
\texttt{exact\_estep=True} (Student-$t$ only) switches to the
covariance-corrected E-step of Section~\ref{sec:robust}.
\texttt{info} contains the final per-period weight vector
\texttt{tau\_v} (and \texttt{tau\_w} when
\texttt{robust\_state=True}); weights much below $1$ flag outliers or
structural breaks.

\paragraph{\texttt{smooth\_slds}.}  Switching dynamics.
\begin{lstlisting}[language=Python]
Xhat, info = sks.smooth_slds(
    Y, F_set, H, Q_set, R, a0, P0,
    A=None, pi0=None, Z=None, B_set=None, D=None,
    max_iter=30, tol=1e-5)
\end{lstlisting}
\texttt{F\_set} and \texttt{Q\_set} are length-$K$ lists of NumPy
arrays; regime-dependent state loadings go in \texttt{B\_set}
(length-$K$ list of $n\times k$ arrays) alongside the common
\texttt{Z} and the regime-invariant \texttt{D}.  The routine
implements the exact variational coordinate updates of
Section~\ref{sec:slds}.  \texttt{info["pi"]} is the $K\times T$
regime-responsibility matrix at convergence;
\texttt{info["Ptt"]} and \texttt{info["Ptt1"]} are the posterior
covariance blocks of $q(X)$.

\paragraph{\texttt{smooth\_constrained}.}  Inequality-constrained smoother.
\begin{lstlisting}[language=Python]
Xhat, info = sks.smooth_constrained(
    Y, F, H, Q, R, a0, P0,
    A_ineq, b_ineq,
    Z=None, B=None, D=None)
\end{lstlisting}
\texttt{A\_ineq} should be a SciPy sparse matrix or a dense array of
shape $(c, (T+1)n)$; \texttt{b\_ineq} is a length-$c$ array.
\texttt{info} contains the active-set mask, approximate KKT
multipliers, and interior-point convergence diagnostics.

\subsection{Demos and tests}

The same eleven demos as the Octave package are provided in
\texttt{python/examples/}, with filenames \texttt{demo\_*.py}
matching the Octave counterparts one-for-one.  Like their Octave
counterparts, the Python demos print numerical diagnostics and
produce \texttt{matplotlib} figures illustrating the same
quantities.  Figures are simultaneously saved to
\texttt{python/examples/figures/} and displayed on screen when a
GUI backend is available; on headless systems the demos fall back
to the non-interactive Agg backend silently.  The three worked
applications of Appendix~\ref{sec:apps} produce figures in the
same way.  Run any demo individually, or use the interactive menu:
\begin{lstlisting}
cd python/examples
python demo_main.py
\end{lstlisting}
Given identical inputs, the Python and Octave/MATLAB core routines
agree to floating-point precision (for example, on the shared
\texttt{demo\_slds} data the smoothed paths agree to $\sim 10^{-16}$
and the regime responsibilities to $\sim 10^{-15}$).  Note, however,
that most demos draw their simulated data from each language's
\emph{native} random-number generator, so their printed numbers differ
across languages while telling the same qualitative story;
\texttt{demo\_slds} is the exception, using a portable generator so
that all environments produce identical output.  The headline numbers
quoted in this note are from the Python runs.  The three worked
applications of
Appendix~\ref{sec:apps} are likewise implemented as
\texttt{app$k$\_*.py} in the \texttt{applications/} subdirectory of
\texttt{python/examples/}.

A small test suite verifies the Python implementation:
\begin{lstlisting}
cd python
python tests/test_all_demos.py
\end{lstlisting}
Five tests are run: simultaneous-vs-RTS agreement to $\sim 10^{-13}$;
selected inversion vs dense inverse; exogenous-consistency
($B = D = 0$ reproduces the no-exogenous solution); EM marginal
log-likelihood is monotone non-decreasing across iterations; and all
eleven demos run without error.

\section{Three worked applications}
\label{sec:apps}

The appendix walks through three end-to-end applications, each of
which is a recognizable workhorse in applied macroeconomic work.  The
purpose is not to introduce new theory --- everything follows from
the machinery developed in Sections \ref{sec:precision-form} through
\ref{sec:constrained} --- but to show, by repeated example, that the
machinery handles substantively different problems with the same tools.

Each application is implemented as a runnable script in both Octave
and Python, in the directories
\texttt{octave/examples/applications/} and
\texttt{python/examples/}\allowbreak\texttt{applications/}, with filenames
\texttt{app$k$\_*.m} and \texttt{app$k$\_*.py} respectively.  Running
them reproduces the numbers reported here.

\subsection{Natural rate of interest}
\label{sec:apps-rstar}

The first application is the natural rate of interest $r^*_t$, the
real short rate consistent with output at its trend.  Following
\citet{LaubachWilliams2003} and \citet{HolstonLaubachWilliams2017}
(``HLW''), we estimate $r^*_t$ jointly with trend output $y^*_t$
and trend growth $g_t$, using a small structural state-space model.

\paragraph{Model.}  Three latent objects evolve as random walks: trend
output $y^*_t$, trend growth $g_t$, and an AR(2) output gap
$\widetilde y_t$.  The natural rate is tied to trend growth through
\begin{equation}
r^*_t = c\,g_t,
\end{equation}
the ``Laubach--Williams identification'' --- a cross-equation
restriction that pins down what would otherwise be an
unidentified latent path.  The state and measurement equations are
\begin{align}
y^*_t &= y^*_{t-1}+g_{t-1}+\varepsilon^{y^*}_t,
   &\varepsilon^{y^*}_t&\sim N(0,\sigma_{y^*}^2),\\
g_t &= g_{t-1}+\varepsilon^g_t,
   &\varepsilon^g_t&\sim N(0,\sigma_g^2),\\
\widetilde y_t &= a_1\widetilde y_{t-1}+a_2\widetilde y_{t-2}
                 +\frac{a_r}{2}\sum_{j=1}^{2}\bigl(r_{t-j}-r^*_{t-j}\bigr)
                 +\varepsilon^{\widetilde y}_t,\\
y_t&=y^*_t+\widetilde y_t,\\
\pi_t - b_\pi\pi_{t-1}-(1-b_\pi)\bar\pi &= b_y\widetilde y_{t-1}+v_t.
\end{align}
The IS curve uses the standard HLW timing: the real-rate \emph{gap}
enters as the average of its first two lags, so with
$r^*_t=c\,g_t$ the equation expands to
\begin{equation}
\widetilde y_t = a_1\widetilde y_{t-1}+a_2\widetilde y_{t-2}
   +\frac{a_r}{2}\bigl(r_{t-1}+r_{t-2}\bigr)
   -\frac{a_r c}{2}\bigl(g_{t-1}+g_{t-2}\bigr)
   +\varepsilon^{\widetilde y}_t :
\end{equation}
the observed policy-rate average enters as an exogenous input, and the
natural-rate average appears as a state-side loading of $-a_r c/2$ on
\emph{each} of $g_{t-1}$ and $g_{t-2}$ --- a single normalization used
consistently in the equations, the transition matrix below, and the
package code.  The inflation equation is implemented in
quasi-differenced form: the left-hand side
$\pi_t-b_\pi\pi_{t-1}-(1-b_\pi)\bar\pi$ is the observable fed to the
smoother (so lagged inflation never needs to live in the state), and
its loading $b_y$ on $\widetilde y_{t-1}$ is what identifies the
Phillips slope.

\paragraph{State-space mapping.}  Augmenting the state to carry the
necessary lags gives
\begin{equation}
x_t=(y^*_t,\ g_t,\ g_{t-1},\ \widetilde y_t,\ \widetilde y_{t-1})^\top\in\mathbb R^5.
\end{equation}
In $x_{t-1}=(y^*_{t-1},\,g_{t-1},\,g_{t-2},\,\widetilde y_{t-1},\,
\widetilde y_{t-2})^\top$, trend growth $g_{t-1}$ is the second
component and $g_{t-2}$ the third, so both lags the IS curve needs are
already in the state.  Reading the equations off row by row, the
transition matrix is
\begin{equation}
F=\begin{pmatrix}
1 & 1 & 0 & 0 & 0\\
0 & 1 & 0 & 0 & 0\\
0 & 1 & 0 & 0 & 0\\
0 & -\tfrac{a_r c}{2} & -\tfrac{a_r c}{2} & a_1 & a_2\\
0 & 0 & 0 & 1 & 0
\end{pmatrix},
\quad B=\begin{pmatrix}0\\0\\0\\a_r\\0\end{pmatrix},
\end{equation}
with exogenous regressor $z_t=\tfrac12(r_{t-1}+r_{t-2})$: the first
row implements $y^*_t=y^*_{t-1}+g_{t-1}$ (the loading sits on the
second component of $x_{t-1}$, i.e.\ $g_{t-1}$, matching the equation's
timing), and the fourth row carries the two $-a_r c/2$ loadings on
$g_{t-1}$ and $g_{t-2}$.  $Q$ is diagonal with \emph{three}
substantive innovation variances ---
$\sigma_{y^*}^2$, $\sigma_g^2$, and $\sigma_{\widetilde y}^2$ ---
plus tiny regularizers on the two deterministic lag-shift components
(rows three and five carry no shock of their own); the inflation
disturbance $v_t$ is a \emph{measurement} shock and lives in $R$, not
$Q$.  The observation matrix has output loading on $y^*$ and
$\widetilde y$ and a quasi-differenced-inflation row that picks up
$b_y\widetilde y_{t-1}$.

\paragraph{Identification.}  Three issues are worth flagging:
\begin{itemize}[leftmargin=2em]
\item The IS-curve slope $a_r$ is famously near zero on
  postwar US data \citep{HolstonLaubachWilliams2017}, which causes a
  pile-up problem on $\sigma_g$ when both are estimated.  The
  simultaneous form does not resolve this --- it is a fundamental
  identification limitation --- but it makes the source visible: the
  EM-estimated likelihood becomes nearly flat in a direction of
  parameter space.
\item Diffuse initialization on $y^*$ and $g$ is required (the
  random-walk trend has no proper stationary distribution).  This is
  the $P_0=\infty$ case treated in Section \ref{sec:diffuse}: the
  prior block of $J$ is dropped.
\item As an illustration of constrained smoothing, we impose
  $r^*_t\ge 0$ on the smoothed path: $-c\,g_t\le 0$ for all $t$,
  plugged into the QP framework of Section~\ref{sec:constrained} with
  one row of $A_{\text{ineq}}$ per period.  One clarification on the
  economics: this is a \emph{maintained restriction on the natural
  rate}, not the zero lower bound --- the ZLB constrains the
  \emph{nominal policy rate} $i_t$, while the natural real rate
  $r^*_t$ can perfectly well be negative in the model (and in the
  simulation below the true $r^*$ sometimes is).  A researcher might
  still maintain $r^*\ge 0$ as a structural prior for a particular
  economy; the point of the exercise is to show how cheaply such a
  restriction is imposed, and (below) what happens when it is wrong.
\end{itemize}

\paragraph{Results on simulated data.}  On 200 quarters of data
generated from the model with $\sigma_g=0.08$ (large enough that
the true $r^*$ visits negative values), the demo reports:

\begin{itemize}[leftmargin=2em]
\item Estimator A (Gaussian smoother, true parameters):
  RMSE on $g_t$ is 0.23; the posterior standard deviation of $r^*$
  at the end of sample is 0.25, roughly 1.5 times its full-sample
  average of 0.17.  This is the well-known result that $r^*$ is more
  precisely identified \emph{inside} the sample than at the latest
  date, where only one-sided information is available.
\item Estimator B (constrained smoother): the unconstrained
  estimator places $r^*$ below zero in 37 of 200 quarters; the
  constrained estimator pins $r^*\ge 0$ to machine precision.
  In this particular Monte Carlo draw the true $r^*$ is itself
  negative in 27 quarters, so the maintained restriction contradicts
  the data-generating process and (mildly) \emph{hurts} RMSE.  This is
  the right diagnostic: an inequality constraint is informative only
  when it reflects truth.  On samples where the true $r^*$ stays
  non-negative (a more plausible empirical setting in some advanced
  economies), the constraint improves accuracy.
\item Estimator C (EM, all parameters unknown):  the EM-estimated $y^*$
  correlates with the true $y^*$ at $\sim 0.99$, with monotone
  marginal log-likelihood across iterations.  Identification of the
  latent dynamics rotation is not perfect --- $F$ in the unrestricted
  EM is not unique, only its action on the data --- but the
  decomposition into trend and cycle is well-recovered.
\end{itemize}

The point of this application is that the same three pieces of
machinery developed earlier in the note --- the basic smoother
(Section~\ref{sec:precision-form}), constrained smoothing
(Section~\ref{sec:constrained}), and EM with exogenous inputs
(Section~\ref{sec:em}) --- combine without modification to deliver a
full HLW-style analysis.  No bespoke "natural-rate code" is needed.

\subsection{Nowcasting current-quarter GDP}
\label{sec:apps-nowcast}

The second application is a mixed-frequency dynamic factor model for
nowcasting current-quarter GDP from monthly indicators.  This is the
workhorse setup at central banks and IMF country desks; the canonical
reference is \citet{MarianoMurasawa2003}, with extensions to large
panels by \citet{GiannoneReichlinSmall2008} and
\citet{BanburaModugno2014}.

\paragraph{Model.}  A single latent monthly factor $f_t$ drives all
indicators:
\begin{align}
f_t&=\phi\,f_{t-1}+e_t,\quad e_t\sim N(0,\sigma_e^2),\\
y_{k,t}&=\lambda_k\,f_t+u_{k,t},\quad u_{k,t}\sim N(0,R_{kk}),\quad k=1,\ldots,K.
\end{align}
Quarterly GDP enters as a third-month aggregate: at the end of each
quarter,
\begin{equation}
\text{GDP}_t=\lambda_g\cdot\tfrac13(f_t+f_{t-1}+f_{t-2})+u^{\text{GDP}}_t,
\end{equation}
and is unobserved (\texttt{NaN}) at all other monthly time stamps.

\paragraph{State augmentation.}  The third-month average requires the
state to carry two lags of the factor:
$x_t=(f_t,f_{t-1},f_{t-2})^\top$, so $n=3$.  The observation matrix has
$K$ rows for the monthly indicators (each loading on the first
component only) and one row for GDP:
\begin{equation}
H=\begin{pmatrix}
\lambda_1 & 0 & 0\\
\vdots & \vdots & \vdots\\
\lambda_K & 0 & 0\\
\lambda_g/3 & \lambda_g/3 & \lambda_g/3
\end{pmatrix}.
\end{equation}

\paragraph{Jagged edge.}  At any real-time vintage, the data panel
has a ``jagged edge'': different indicators have different
publication lags (industrial production typically lags by 1 month,
employment by 1 month, retail sales by 1-2 months, GDP by 1-2
quarters), so the last several observations are missing for some
series and present for others.  In the simultaneous formulation this
is just a pattern of NaN entries; the smoother handles them by
dropping rows of $H'R^{-1}H$ at each affected $t$.  No bespoke
"missing-data" code is needed.

\paragraph{Real-time nowcast experiment.}  At the end of each
of the three months of a quarter, we re-run the smoother on the
data available at that vintage and read off the current-quarter
nowcast as the smoothed average of the factor over the three
months.  The posterior variance of this average is
\begin{equation}
\Var(\widehat{\text{GDP}})=\lambda_g^2\cdot\tfrac19\,\mathbf 1^\top\Sigma_3\,\mathbf 1,
\end{equation}
where $\Sigma_3$ is the $3\times 3$ posterior covariance of the
three factor values over the quarter --- and here the state
augmentation pays a small dividend.  Because the augmented state at
the end-of-quarter month $t_3$ is
$x_{t_3}=(f_{t_3},f_{t_3-1},f_{t_3-2})^\top$, $\Sigma_3$ is the
\emph{single} diagonal block $\Sigma_{t_3 t_3}$ delivered by the
Takahashi recursion of Section~\ref{sec:selinv}; its off-diagonal
\emph{elements} already contain all three cross-month covariances,
including the lag-two term $\Cov(f_{t_3},f_{t_3-2}\mid Y)$.  No
off-diagonal time blocks of $J^{-1}$ are needed for this formula
(assembling $\Sigma_3$ from the scalar factor at three separate times
would require them, and dropping the lag-two cross term --- which is
positive for a persistent factor --- would understate the nowcast
variance).

\paragraph{Results on simulated data.}  On 10 years of monthly
data with $K=5$ indicators (release lags $[0,1,1,0,0]$) and one
quarterly GDP series, the nowcast standard deviation shrinks from
0.36 with one month of data to 0.21 with two months to 0.12 at the
end of the quarter.  The factor itself is recovered with 0.96
correlation against truth despite the missing-data pattern.
Nowcast point estimates fall within one standard deviation of the
true value at every vintage.

The point of this application is the modularity claim from
Section~\ref{sec:missing}: jagged edges, mixed frequencies, and
publication lags are all the same thing --- patterns of missing
observations --- and the simultaneous smoother handles them with no
special-case code.  The posterior covariance from selected inversion
gives the calibrated nowcast standard error as a free byproduct.

\subsection{Semi-structural three-equation New Keynesian model}
\label{sec:apps-nk}

The third application is the workhorse three-equation New Keynesian
model used for semi-structural projections in policy institutions
(see \citet{BergKaramLaxton2006} for the IMF reference
implementation).  Three equations link the output gap, inflation,
and the policy rate:
\begin{align}
\widetilde y_t &= \alpha_1\widetilde y_{t-1}+\alpha_2(i_{t-1}-\pi_{t-1}-\bar r)+\varepsilon^{y}_t,\\
\pi_t &= \beta_1\pi_{t-1}+(1-\beta_1)\bar\pi+\beta_2\widetilde y_{t-1}+\varepsilon^{\pi}_t,\\
i_t &= \rho_i i_{t-1}+(1-\rho_i)\bigl[\bar r+\bar\pi+\phi_\pi(\pi_{t-1}-\bar\pi)+\phi_y\widetilde y_{t-1}\bigr]+\varepsilon^{i}_t.
\end{align}
The structural slopes $(\beta_2,\alpha_2,\phi_\pi,\phi_y)$ are the
empirically interesting parameters; the structural shocks
$(\varepsilon^y,\varepsilon^\pi,\varepsilon^i)$ are demand,
cost-push, and monetary shocks, the bread and butter of policy
narratives.

\paragraph{State-space mapping.}  The state $x_t=(\widetilde y_t,\pi_t,i_t)$
contains everything an observer needs.  All forward
expectations have been replaced by their backward-looking reduced
form; the system is purely autoregressive with three innovations.
The $F$ matrix is
\begin{equation}
F=\begin{pmatrix}
\alpha_1 & -\alpha_2 & \alpha_2\\
\beta_2 & \beta_1 & 0\\
(1-\rho_i)\phi_y & (1-\rho_i)\phi_\pi & \rho_i
\end{pmatrix},
\quad
B=\begin{pmatrix}
-\alpha_2\bar r\\
(1-\beta_1)\bar\pi\\
(1-\rho_i)(\bar r+\bar\pi-\phi_\pi\bar\pi)
\end{pmatrix},
\end{equation}
with the exogenous regressor $z_t=1$ carrying the constants.
Observation is direct: $y_t=x_t+v_t$, with measurement error.

\paragraph{Shock decomposition.}  Given smoothed states $\widehat X$,
the smoothed structural shocks fall out:
\begin{equation}
\widehat\varepsilon_t=\widehat x_t-F\widehat x_{t-1}-B z_t.
\end{equation}
These are the central-bank-narrative quantities: "how much of last
quarter's output gap was demand vs monetary?"  The simultaneous
smoother delivers them in one sparse solve.

\paragraph{Results on simulated data.}  On 200 quarters of simulated
data with conventional calibration ($\beta_2=0.15$, $\rho_i=0.80$,
$\phi_\pi=1.50$, $\phi_y=0.50$):

\begin{itemize}[leftmargin=2em]
\item The Gaussian smoother with true parameters recovers all three
  structural shocks with correlations of $0.90$, $0.99$, and $0.97$
  against truth, respectively (demand, cost-push, monetary).
\item EM with all parameters unknown converges in $\sim 120$
  iterations with monotone marginal log-likelihood.  The Phillips
  slope $\beta_2$ (i.e.\ $\kappa$, the canonical hard parameter) is
  recovered at $0.17$ vs.\ a truth of $0.15$.  The Taylor-rule
  coefficients are noticeably less precise ($\phi_\pi$ at $1.87$
  vs.\ truth $1.50$), a well-known consequence of the limited
  variation in policy-rate movements available in the simulated
  sample: with the smoothing parameter $\rho_i=0.80$, the
  systematic part of the rule contributes little independent
  variation, and the likelihood is correspondingly flat in
  $(\phi_\pi,\phi_y)$.
\end{itemize}

No DSGE-specific estimation routine is involved.  The state-space
formulation plus the standard EM algorithm of Section~\ref{sec:em}
delivers structural identification, shock decomposition, and
posterior bands on the latent gaps on the same sparse $J$ that the
basic smoother used.

\subsection{What the three applications have in common}

The point of the three together is conceptual unity.  Each
application defines $(F,H,Q,R,B,D)$ for a substantively different
economic problem.  Each produces a substantively different output:
$r^*$ at the latest date, current-quarter GDP, structural shock
decomposition.  Each exercises a different optional capability of
the framework: constrained smoothing, missing data, EM on a
multi-equation system.  But the core machinery is identical.  The
work-horse object in every case is the same sparse block-tridiagonal
linear system $JX=h$, solved once or iterated as part of EM, with
the same selected-inversion machinery for posterior covariances.
This is what is meant by ``the right idiom'' for state-space modelling
in macroeconometrics.

\begin{thebibliography}{Lange et~al.(1989)}

\bibitem[Andrews and Mallows(1974)]{AndrewsMallows1974}
Andrews, D.\ F., and C.\ L.\ Mallows (1974).
``Scale mixtures of normal distributions.''
\emph{Journal of the Royal Statistical Society, Series B},
36(1):99--102.

\bibitem[Banbura and Modugno(2014)]{BanburaModugno2014}
Banbura, M., and M.\ Modugno (2014).
``Maximum likelihood estimation of factor models on datasets with
arbitrary pattern of missing data.''
\emph{Journal of Applied Econometrics}, 29(1):133--160.

\bibitem[Berg, Karam, and Laxton(2006)]{BergKaramLaxton2006}
Berg, A., P.\ Karam, and D.\ Laxton (2006).
``Practical model-based monetary policy analysis: a how-to guide.''
\emph{IMF Working Paper} WP/06/80.

\bibitem[Bierman(1977)]{Bierman1977}
Bierman, G.\ J.\ (1977).
\emph{Factorization Methods for Discrete Sequential Estimation}.
Academic Press, New York.

\bibitem[Botev(2017)]{Botev2017}
Botev, Z.\ I.\ (2017).
``The normal law under linear restrictions: simulation and estimation
via minimax tilting.''
\emph{Journal of the Royal Statistical Society, Series B},
79(1):125--148.

\bibitem[Carpenter et~al.(2017)]{Carpenter2017}
Carpenter, B., A.\ Gelman, M.\ D.\ Hoffman, D.\ Lee, B.\ Goodrich,
M.\ Betancourt, M.\ Brubaker, J.\ Guo, P.\ Li, and A.\ Riddell (2017).
``Stan: a probabilistic programming language.''
\emph{Journal of Statistical Software}, 76(1):1--32.

\bibitem[Chan and Jeliazkov(2009)]{ChanJeliazkov2009}
Chan, J.\ C.\ C., and I.\ Jeliazkov (2009).
``Efficient simulation and integrated likelihood estimation in state
space models.''
\emph{International Journal of Mathematical Modelling and Numerical
Optimisation}, 1:101--120.

\bibitem[Davis(2006)]{Davis2006}
Davis, T.\ A.\ (2006).
\emph{Direct Methods for Sparse Linear Systems}.
SIAM.

\bibitem[de Jong(1991)]{deJong1989}
de Jong, P.\ (1991).
``The diffuse Kalman filter.''
\emph{Annals of Statistics}, 19(2):1073--1083.

\bibitem[Durbin and Koopman(2012)]{DurbinKoopman2012}
Durbin, J., and S.\ J.\ Koopman (2012).
\emph{Time Series Analysis by State Space Methods}, 2nd ed.
Oxford University Press.

\bibitem[Erisman and Tinney(1975)]{ErismanTinney1975}
Erisman, A.\ M., and W.\ F.\ Tinney (1975).
``On computing certain elements of the inverse of a sparse matrix.''
\emph{Communications of the ACM}, 18(3):177--179.

\bibitem[Ghahramani and Hinton(2000)]{GhahramaniHinton2000}
Ghahramani, Z., and G.\ E.\ Hinton (2000).
``Variational learning for switching state-space models.''
\emph{Neural Computation}, 12(4):831--864.

\bibitem[Giannone, Reichlin, and Small(2008)]{GiannoneReichlinSmall2008}
Giannone, D., L.\ Reichlin, and D.\ Small (2008).
``Nowcasting: the real-time informational content of macroeconomic
data.''
\emph{Journal of Monetary Economics}, 55(4):665--676.

\bibitem[Giannone, Lenza, and Primiceri(2015)]{GiannoneLenzaPrimiceri2015}
Giannone, D., M.\ Lenza, and G.\ E.\ Primiceri (2015).
``Prior selection for vector autoregressions.''
\emph{Review of Economics and Statistics}, 97(2):436--451.

\bibitem[Holland and Welsch(1977)]{HollandWelsch1977}
Holland, P.\ W., and R.\ E.\ Welsch (1977).
``Robust regression using iteratively reweighted least-squares.''
\emph{Communications in Statistics: Theory and Methods}, 6(9):813--827.

\bibitem[Holston, Laubach, and Williams(2017)]{HolstonLaubachWilliams2017}
Holston, K., T.\ Laubach, and J.\ C.\ Williams (2017).
``Measuring the natural rate of interest: international trends and
determinants.''
\emph{Journal of International Economics}, 108(S1):S59--S75.

\bibitem[Kalman(1960)]{Kalman1960}
Kalman, R.\ E.\ (1960).
``A new approach to linear filtering and prediction problems.''
\emph{Journal of Basic Engineering}, 82(D):35--45.

\bibitem[Koopman and Durbin(2003)]{KoopmanDurbin2003}
Koopman, S.\ J., and J.\ Durbin (2003).
``Filtering and smoothing of state vector for diffuse state space
models.''
\emph{Journal of Time Series Analysis}, 24:85--98.

\bibitem[Lange et~al.(1989)]{LangeLittleTaylor1989}
Lange, K.\ L., R.\ J.\ A.\ Little, and J.\ M.\ G.\ Taylor (1989).
``Robust statistical modeling using the $t$-distribution.''
\emph{Journal of the American Statistical Association},
84(408):881--896.

\bibitem[Laubach and Williams(2003)]{LaubachWilliams2003}
Laubach, T., and J.\ C.\ Williams (2003).
``Measuring the natural rate of interest.''
\emph{Review of Economics and Statistics}, 85(4):1063--1070.

\bibitem[Mariano and Murasawa(2003)]{MarianoMurasawa2003}
Mariano, R.\ S., and Y.\ Murasawa (2003).
``A new coincident index of business cycles based on monthly and
quarterly series.''
\emph{Journal of Applied Econometrics}, 18(4):427--443.

\bibitem[McCausland et~al.(2011)]{McCauslandMillerPelletier2011}
McCausland, W.\ J., S.\ Miller, and D.\ Pelletier (2011).
``Simulation smoothing for state-space models: A computational
efficiency analysis.''
\emph{Computational Statistics \& Data Analysis}, 55(1):199--212.

\bibitem[Rauch et~al.(1965)]{RTS1965}
Rauch, H.\ E., F.\ Tung, and C.\ T.\ Striebel (1965).
``Maximum likelihood estimates of linear dynamic systems.''
\emph{AIAA Journal}, 3(8):1445--1450.

\bibitem[Roweis and Ghahramani(1999)]{RoweisGhahramani1999}
Roweis, S., and Z.\ Ghahramani (1999).
``A unifying review of linear Gaussian models.''
\emph{Neural Computation}, 11(2):305--345.

\bibitem[Rue and Held(2005)]{RueHeld2005}
Rue, H., and L.\ Held (2005).
\emph{Gaussian Markov Random Fields: Theory and Applications}.
Chapman \& Hall/CRC.

\bibitem[Rue, Martino, and Chopin(2009)]{RueMartinoChopin2009}
Rue, H., S.\ Martino, and N.\ Chopin (2009).
``Approximate Bayesian inference for latent Gaussian models by using
integrated nested Laplace approximations.''
\emph{Journal of the Royal Statistical Society, Series B},
71(2):319--392.

\bibitem[Shumway and Stoffer(1982)]{ShumwayStoffer1982}
Shumway, R.\ H., and D.\ S.\ Stoffer (1982).
``An approach to time series smoothing and forecasting using the EM
algorithm.''
\emph{Journal of Time Series Analysis}, 3(4):253--264.

\bibitem[Takahashi et~al.(1973)]{TakahashiFaganChen1973}
Takahashi, K., J.\ Fagan, and M.\ Chen (1973).
``Formation of a sparse bus impedance matrix and its application to
short circuit study.''
\emph{Proc.\ 8th PICA Conference}, pp.\ 63--69.  Minneapolis.

\bibitem[Watson and Engle(1983)]{Watson1983}
Watson, M.\ W., and R.\ F.\ Engle (1983).
``Alternative algorithms for the estimation of dynamic factor,
MIMIC and varying coefficient regression models.''
\emph{Journal of Econometrics}, 23(3):385--400.

\end{thebibliography}

\end{document}
